% Recovered from the previously completed PDF; no new mathematical content.
% Prose is editable; formula boxes retain original symbol positions.
% Build from repository root: xelatex -no-shell-escape -halt-on-error -output-directory=... print/fudan/NAME.tex
\documentclass[11pt,a4paper]{article}
\usepackage[margin=23mm,headheight=14pt,headsep=8.5mm,footskip=10mm]{geometry}
\usepackage{amsmath,amssymb}
\DeclareMathAlphabet{\mathsf}{OT1}{cmss}{m}{n}
\usepackage[no-math]{fontspec}
\usepackage{xeCJK}
\setmainfont{Latin Modern Roman}
\setCJKmainfont{FandolSong-Regular.otf}[Path=assets/fonts/fandol/,BoldFont=FandolSong-Bold.otf,ItalicFont=FandolSong-Regular.otf,Script=CustomDefault]
\xeCJKsetup{PunctStyle=quanjiao,CheckSingle=true}
\usepackage{setspace}\setstretch{1.12}
\setlength{\parindent}{2em}\setlength{\parskip}{3pt}
\emergencystretch=6em
\hfuzz=9pt
\usepackage{titlesec}
\titleformat{\section}{\centering\normalfont\Large\bfseries}{}{0pt}{}
\titlespacing*{\section}{0pt}{1.8\baselineskip}{1.0\baselineskip}
\usepackage{fancyhdr}\pagestyle{fancy}\fancyhf{}
\fancyhead[L]{\small 数学手记·高等代数}\fancyhead[R]{\small 18A08–18A14}
\fancyfoot[C]{\thepage}\renewcommand{\headrulewidth}{0.4pt}
\usepackage[hidelinks,unicode]{hyperref}
\hypersetup{pdftitle={18A08–18A14：题目、解答与小结},pdfauthor={Vincent Arqain}}
\newcommand{\RB}[4]{\raisebox{0pt}[#2bp][#3bp]{\hbox to #1bp{#4\hss}}}
\newcommand{\RG}[5]{\hbox to 0pt{\kern#1bp\raisebox{#2bp}[0pt][0pt]{\fontsize{#3bp}{#3bp}\csname font#4\endcsname\char#5}\hss}}
\newcommand{\RT}[4]{\hbox to 0pt{\kern#1bp\raisebox{#2bp}[0pt][0pt]{\fontsize{#3bp}{#3bp}\selectfont#4}\hss}}
\newcommand{\RS}[4]{\hbox to 0pt{\kern#1bp\raisebox{#2bp}[0pt][0pt]{\fontsize{#3bp}{#3bp}\selectfont\ensuremath{#4}}\hss}}
\newcommand{\RR}[4]{\hbox to 0pt{\kern#1bp\raisebox{#2bp}[0pt][0pt]{\rule{#3bp}{#4bp}}\hss}}
\newcommand{\fontmi}{\usefont{OML}{cmm}{m}{it}}
\newcommand{\fontsy}{\usefont{OMS}{cmsy}{m}{n}}
\newcommand{\fontex}{\usefont{OMX}{cmex}{m}{n}}
\newcommand{\fontrm}{\usefont{OT1}{cmr}{m}{n}}
\newcommand{\fontbf}{\usefont{OT1}{cmr}{bx}{n}}
\newcommand{\fontam}{\usefont{U}{msa}{m}{n}}
\newcommand{\fontbb}{\usefont{U}{msb}{m}{n}}
\widowpenalty=10000\clubpenalty=10000\raggedbottom
\begin{document}
\begin{center}{\LARGE\bfseries 18A08–18A14\par}\vspace{8pt}复旦大学高等代数习题集·题目与解答\end{center}
% Original page 1.
\section*{2018A08}
设\allowbreak{}\hskip.1em \RB{21.773}{8.316}{2.600}{%
\RG{0.000}{0.005}{10.9091}{mi}{65}\RG{8.182}{0.005}{10.9091}{rm}{44}\RG{13.026}{0.005}{10.9091}{mi}{66}%
}\hskip.1em\allowbreak{} 分别为\allowbreak{}\hskip.1em \RB{29.638}{5.861}{0.615}{%
\RG{0.000}{0.005}{10.9091}{mi}{109}\RG{12.000}{0.005}{10.9091}{sy}{2}\RG{22.909}{0.005}{10.9091}{mi}{110}%
}\hskip.1em\allowbreak{} 和\allowbreak{}\hskip.1em \RB{29.628}{5.861}{0.615}{%
\RG{0.000}{0.005}{10.9091}{mi}{110}\RG{8.967}{0.005}{10.9091}{sy}{2}\RG{19.877}{0.005}{10.9091}{mi}{109}%
}\hskip.1em\allowbreak{} 矩阵，\allowbreak{}\hskip.1em \RB{8.791}{8.196}{0.735}{%
\RG{0.000}{0.005}{10.9091}{mi}{67}%
}\hskip.1em\allowbreak{} 为\allowbreak{}\hskip.1em \RB{6.729}{5.327}{0.615}{%
\RG{0.000}{0.005}{10.9091}{mi}{110}%
}\hskip.1em\allowbreak{} 阶非异阵，满足

\[
\RB{80.857}{8.657}{3.208}{%
\RG{0.000}{-0.003}{10.9091}{mi}{65}\RG{8.182}{-0.003}{10.9091}{rm}{40}\RG{12.426}{-0.003}{10.9091}{mi}{67}\RG{23.444}{-0.003}{10.9091}{rm}{43}\RG{34.353}{-0.003}{10.9091}{mi}{66}\RG{42.906}{-0.003}{10.9091}{mi}{65}\RG{51.088}{-0.003}{10.9091}{rm}{41}\RG{58.364}{-0.003}{10.9091}{rm}{61}\RG{69.884}{-0.003}{10.9091}{mi}{79}\RG{78.263}{-0.003}{10.9091}{rm}{46}%
}
\]
证明：线性方程组\allowbreak{}\hskip.1em \RB{34.482}{8.310}{0.741}{%
\RG{0.000}{-0.001}{10.9091}{mi}{65}\RG{8.182}{-0.001}{10.9091}{mi}{120}\RG{17.444}{-0.001}{10.9091}{rm}{61}\RG{28.964}{-0.001}{10.9091}{rm}{48}%
}\hskip.1em\allowbreak{} 的通解为\allowbreak{}\hskip.1em \RB{54.217}{8.659}{3.206}{%
\RG{0.000}{-0.001}{10.9091}{rm}{40}\RG{4.244}{-0.001}{10.9091}{mi}{67}\RG{15.262}{-0.001}{10.9091}{rm}{43}\RG{26.171}{-0.001}{10.9091}{mi}{66}\RG{34.724}{-0.001}{10.9091}{mi}{65}\RG{42.906}{-0.001}{10.9091}{rm}{41}\RG{47.150}{-0.001}{10.9091}{mi}{11}%
}\hskip.1em\allowbreak{} ，其中\allowbreak{}\hskip.1em \RB{7.067}{5.321}{0.621}{%
\RG{0.000}{-0.001}{10.9091}{mi}{11}%
}\hskip.1em\allowbreak{} 为任意的\allowbreak{}\hskip.1em \RB{6.729}{5.321}{0.621}{%
\RG{0.000}{-0.001}{10.9091}{mi}{110}%
}\hskip.1em\allowbreak{} 维列向量。证明：记\allowbreak{}\hskip.1em \RB{62.475}{8.306}{1.410}{%
\RG{0.000}{-0.005}{10.9091}{mi}{68}\RG{12.098}{-0.005}{10.9091}{rm}{61}\RG{23.618}{-0.005}{10.9091}{mi}{67}\RG{34.637}{-0.005}{10.9091}{rm}{43}\RG{45.557}{-0.005}{10.9091}{mi}{66}\RG{54.110}{-0.005}{10.9091}{mi}{65}%
}\hskip.1em\allowbreak{} 。由\allowbreak{}\hskip.1em \RB{40.373}{8.306}{0.745}{%
\RG{0.000}{-0.005}{10.9091}{mi}{65}\RG{8.182}{-0.005}{10.9091}{mi}{68}\RG{20.280}{-0.005}{10.9091}{rm}{61}\RG{31.800}{-0.005}{10.9091}{mi}{79}%
}\hskip.1em\allowbreak{} ，对任意\allowbreak{}\hskip.1em \RB{32.816}{7.946}{0.974}{%
\RG{0.000}{-0.005}{10.9091}{mi}{11}\RG{10.004}{-0.005}{10.9091}{sy}{50}\RG{20.313}{-0.005}{10.9091}{bb}{75}\RG{27.291}{3.955}{7.9701}{mi}{110}%
}\hskip.1em\allowbreak{} ，有

\[
\RB{55.439}{8.657}{3.208}{%
\RG{0.000}{-0.003}{10.9091}{mi}{65}\RG{8.182}{-0.003}{10.9091}{rm}{40}\RG{12.426}{-0.003}{10.9091}{mi}{68}\RG{21.502}{-0.003}{10.9091}{mi}{11}\RG{28.484}{-0.003}{10.9091}{rm}{41}\RG{35.749}{-0.003}{10.9091}{rm}{61}\RG{47.270}{-0.003}{10.9091}{rm}{48}\RG{52.724}{-0.003}{10.9091}{rm}{44}%
}
\]
故Im\allowbreak{}\hskip.1em \RB{20.158}{7.950}{1.886}{%
\RG{0.000}{-0.001}{10.9091}{mi}{68}\RG{12.098}{-0.001}{10.9091}{sy}{18}%
}\hskip.1em\allowbreak{} Ker\allowbreak{}\hskip.1em \RB{8.366}{8.310}{0.501}{%
\RG{0.000}{-0.001}{10.9091}{mi}{65}%
}\hskip.1em\allowbreak{} 。

另一方面，由\allowbreak{}\hskip.1em \RB{62.475}{8.316}{0.735}{%
\RG{0.000}{0.005}{10.9091}{mi}{67}\RG{11.618}{0.005}{10.9091}{rm}{61}\RG{23.138}{0.005}{10.9091}{mi}{68}\RG{34.637}{0.005}{10.9091}{sy}{0}\RG{45.557}{0.005}{10.9091}{mi}{66}\RG{54.110}{0.005}{10.9091}{mi}{65}%
}\hskip.1em\allowbreak{} 及\allowbreak{}\hskip.1em \RB{8.791}{8.196}{0.735}{%
\RG{0.000}{0.005}{10.9091}{mi}{67}%
}\hskip.1em\allowbreak{} 非异，得到

\[
\RB{121.075}{9.290}{23.453}{%
\RG{0.000}{-0.003}{10.9091}{mi}{110}\RG{9.567}{-0.003}{10.9091}{rm}{61}\RT{21.093}{-0.003}{10.9091}{r}\RG{25.370}{-0.003}{10.9091}{rm}{40}\RG{29.614}{-0.003}{10.9091}{mi}{67}\RG{38.199}{-0.003}{10.9091}{rm}{41}\RG{45.476}{-0.003}{10.9091}{sy}{20}\RT{57.002}{-0.003}{10.9091}{r}\RG{61.279}{-0.003}{10.9091}{rm}{40}\RG{65.522}{-0.003}{10.9091}{mi}{68}\RG{74.599}{-0.003}{10.9091}{rm}{41}\RG{81.264}{-0.003}{10.9091}{rm}{43}\RT{92.179}{-0.003}{10.9091}{r}\RG{96.455}{-0.003}{10.9091}{rm}{40}\RG{100.698}{-0.003}{10.9091}{mi}{66}\RG{109.251}{-0.003}{10.9091}{mi}{65}\RG{117.433}{-0.003}{10.9091}{rm}{41}\RG{45.484}{-20.247}{10.9091}{sy}{20}\RT{57.002}{-20.247}{10.9091}{r}\RG{61.279}{-20.247}{10.9091}{rm}{40}\RG{65.522}{-20.247}{10.9091}{mi}{68}\RG{74.599}{-20.247}{10.9091}{rm}{41}\RG{81.264}{-20.247}{10.9091}{rm}{43}\RT{92.179}{-20.247}{10.9091}{r}\RG{96.455}{-20.247}{10.9091}{rm}{40}\RG{100.698}{-20.247}{10.9091}{mi}{65}\RG{108.880}{-20.247}{10.9091}{rm}{41}\RG{113.124}{-20.247}{10.9091}{rm}{46}%
}
\]
这里，第一个不等式来自\allowbreak{}\hskip.1em \RB{39.675}{8.316}{0.500}{%
\RG{0.000}{0.005}{10.9091}{mi}{68}\RG{11.673}{0.005}{10.9091}{sy}{0}\RG{22.757}{0.005}{10.9091}{mi}{66}\RG{31.309}{0.005}{10.9091}{mi}{65}%
}\hskip.1em\allowbreak{} 的列空间包含于\allowbreak{}\hskip.1em \RB{9.271}{7.956}{0.500}{%
\RG{0.000}{0.005}{10.9091}{mi}{68}%
}\hskip.1em\allowbreak{} 与\allowbreak{}\hskip.1em \RB{16.907}{8.316}{0.500}{%
\RG{0.000}{0.005}{10.9091}{mi}{66}\RG{8.542}{0.005}{10.9091}{mi}{65}%
}\hskip.1em\allowbreak{} 的列空间之和；第二个不等式来自\allowbreak{}\hskip.1em \RB{8.747}{7.947}{0.504}{%
\RG{0.000}{-0.004}{10.9091}{mi}{66}%
}\hskip.1em\allowbreak{} 将\allowbreak{}\hskip.1em \RB{8.366}{8.307}{0.504}{%
\RG{0.000}{-0.004}{10.9091}{mi}{65}%
}\hskip.1em\allowbreak{} 的列空间映入\allowbreak{}\hskip.1em \RB{16.907}{8.307}{0.504}{%
\RG{0.000}{-0.004}{10.9091}{mi}{66}\RG{8.542}{-0.004}{10.9091}{mi}{65}%
}\hskip.1em\allowbreak{} 的列空间，线性映射不增加维数。因此

\[
\RB{200.443}{9.291}{3.207}{%
\RT{0.000}{-0.002}{10.9091}{d}\RT{6.065}{-0.002}{10.9091}{i}\RT{9.098}{-0.002}{10.9091}{m}\RT{19.997}{-0.002}{10.9091}{I}\RT{23.935}{-0.002}{10.9091}{m}\RG{34.847}{-0.002}{10.9091}{mi}{68}\RG{46.946}{-0.002}{10.9091}{rm}{61}\RT{58.472}{-0.002}{10.9091}{r}\RG{62.749}{-0.002}{10.9091}{rm}{40}\RG{66.992}{-0.002}{10.9091}{mi}{68}\RG{76.069}{-0.002}{10.9091}{rm}{41}\RG{83.334}{-0.002}{10.9091}{sy}{21}\RG{94.854}{-0.002}{10.9091}{mi}{110}\RG{103.822}{-0.002}{10.9091}{sy}{0}\RT{114.743}{-0.002}{10.9091}{r}\RG{119.019}{-0.002}{10.9091}{rm}{40}\RG{123.263}{-0.002}{10.9091}{mi}{65}\RG{131.445}{-0.002}{10.9091}{rm}{41}\RG{138.710}{-0.002}{10.9091}{rm}{61}\RT{150.237}{-0.002}{10.9091}{d}\RT{156.303}{-0.002}{10.9091}{i}\RT{159.335}{-0.002}{10.9091}{m}\RT{170.234}{-0.002}{10.9091}{K}\RT{178.721}{-0.002}{10.9091}{e}\RT{183.565}{-0.002}{10.9091}{r}\RG{189.667}{-0.002}{10.9091}{mi}{65}\RG{197.848}{-0.002}{10.9091}{rm}{46}%
}
\]
最后一个等式是秩与零度公式：把Ker\allowbreak{}\hskip.1em \RB{8.366}{8.311}{0.500}{%
\RG{0.000}{0.000}{10.9091}{mi}{65}%
}\hskip.1em\allowbreak{} 的一组基扩充为\allowbreak{}\hskip.1em \RB{12.496}{7.951}{0.500}{%
\RG{0.000}{0.000}{10.9091}{bb}{75}\RG{6.971}{3.961}{7.9701}{mi}{110}%
}\hskip.1em\allowbreak{} 的一组基，其余基向量在\allowbreak{}\hskip.1em \RB{8.366}{8.311}{0.500}{%
\RG{0.000}{0.000}{10.9091}{mi}{65}%
}\hskip.1em\allowbreak{} 下的像构成Im\allowbreak{}\hskip.1em \RB{8.366}{8.312}{0.500}{%
\RG{0.000}{0.001}{10.9091}{mi}{65}%
}\hskip.1em\allowbreak{} 的一组基。

结合Im\allowbreak{}\hskip.1em \RB{20.158}{7.948}{1.888}{%
\RG{0.000}{-0.003}{10.9091}{mi}{68}\RG{12.098}{-0.003}{10.9091}{sy}{18}%
}\hskip.1em\allowbreak{} Ker\allowbreak{}\hskip.1em \RB{8.366}{8.308}{0.503}{%
\RG{0.000}{-0.003}{10.9091}{mi}{65}%
}\hskip.1em\allowbreak{} ，可得

\[
\RB{104.930}{9.291}{3.207}{%
\RT{0.000}{-0.002}{10.9091}{I}\RT{3.938}{-0.002}{10.9091}{m}\RG{13.025}{-0.002}{10.9091}{rm}{40}\RG{17.269}{-0.002}{10.9091}{mi}{67}\RG{28.287}{-0.002}{10.9091}{rm}{43}\RG{39.196}{-0.002}{10.9091}{mi}{66}\RG{47.749}{-0.002}{10.9091}{mi}{65}\RG{55.931}{-0.002}{10.9091}{rm}{41}\RG{63.207}{-0.002}{10.9091}{rm}{61}\RT{74.729}{-0.002}{10.9091}{K}\RT{83.216}{-0.002}{10.9091}{e}\RT{88.060}{-0.002}{10.9091}{r}\RG{94.154}{-0.002}{10.9091}{mi}{65}\RG{102.336}{-0.002}{10.9091}{rm}{46}%
}
\]
于是每个解都可写成\allowbreak{}\hskip.1em \RB{54.217}{8.664}{3.201}{%
\RG{0.000}{0.004}{10.9091}{rm}{40}\RG{4.244}{0.004}{10.9091}{mi}{67}\RG{15.262}{0.004}{10.9091}{rm}{43}\RG{26.171}{0.004}{10.9091}{mi}{66}\RG{34.724}{0.004}{10.9091}{mi}{65}\RG{42.906}{0.004}{10.9091}{rm}{41}\RG{47.150}{0.004}{10.9091}{mi}{11}%
}\hskip.1em\allowbreak{} ，而每个这样的向量也都是解，故所给表达式就是通解：

小结：先由乘积为零得到像空间包含于核，再利用\allowbreak{}\hskip.1em \RB{8.791}{8.192}{0.739}{%
\RG{0.000}{0.001}{10.9091}{mi}{67}%
}\hskip.1em\allowbreak{} 非异获得反向的维数估计。

\[
\RB{278.451}{8.393}{2.237}{%
\RT{0.000}{0.002}{7.9701}{题}\RT{8.966}{0.002}{7.9701}{目}\RT{17.933}{0.002}{7.9701}{来}\RT{26.899}{0.002}{7.9701}{源}\RT{35.866}{0.002}{7.9701}{：}\RT{44.832}{0.002}{7.9701}{谢}\RT{53.799}{0.002}{7.9701}{启}\RT{62.765}{0.002}{7.9701}{鸿}\RT{71.732}{0.002}{7.9701}{《}\RT{80.698}{0.002}{7.9701}{复}\RT{89.665}{0.002}{7.9701}{旦}\RT{98.631}{0.002}{7.9701}{大}\RT{107.598}{0.002}{7.9701}{学}\RT{116.564}{0.002}{7.9701}{高}\RT{125.531}{0.002}{7.9701}{等}\RT{134.497}{0.002}{7.9701}{代}\RT{143.464}{0.002}{7.9701}{数}\RT{152.430}{0.002}{7.9701}{习}\RT{161.397}{0.002}{7.9701}{题}\RT{170.363}{0.002}{7.9701}{集}\RT{179.329}{0.002}{7.9701}{（}\RT{188.296}{0.002}{7.9701}{2}\RT{192.779}{0.002}{7.9701}{0}\RT{197.262}{0.002}{7.9701}{2}\RT{201.745}{0.002}{7.9701}{3}\RT{209.214}{0.002}{7.9701}{版}\RT{218.180}{0.002}{7.9701}{）}\RT{222.672}{0.002}{7.9701}{》}\RT{227.156}{0.002}{7.9701}{，}\RT{236.122}{0.002}{7.9701}{第}\RT{248.066}{0.002}{7.9701}{6}\RT{252.549}{0.002}{7.9701}{8}\RT{260.018}{0.002}{7.9701}{页}\RT{268.985}{0.002}{7.9701}{。}%
}
\]
% Original page 2.
\section*{2018A09}
设\allowbreak{}\hskip.1em \RB{7.536}{8.196}{0.735}{%
\RG{0.000}{0.005}{10.9091}{mi}{83}%
}\hskip.1em\allowbreak{} 是线性空间\allowbreak{}\hskip.1em \RB{8.889}{7.956}{0.735}{%
\RG{0.000}{0.005}{10.9091}{mi}{86}%
}\hskip.1em\allowbreak{} 中的向量族，并且至少包含一个非零向量。证明：\allowbreak{}\hskip.1em \RB{7.536}{8.196}{0.735}{%
\RG{0.000}{0.005}{10.9091}{mi}{83}%
}\hskip.1em\allowbreak{} 存在极大无关组的充分必要条件是\allowbreak{}\hskip.1em \RB{7.536}{8.187}{0.744}{%
\RG{0.000}{-0.004}{10.9091}{mi}{83}%
}\hskip.1em\allowbreak{} 张成的子空间\allowbreak{}\hskip.1em \RB{22.646}{8.656}{3.209}{%
\RG{0.000}{-0.004}{10.9091}{mi}{76}\RG{7.429}{-0.004}{10.9091}{rm}{40}\RG{11.673}{-0.004}{10.9091}{mi}{83}\RG{19.004}{-0.004}{10.9091}{rm}{41}%
}\hskip.1em\allowbreak{} 是\allowbreak{}\hskip.1em \RB{8.889}{7.947}{0.744}{%
\RG{0.000}{-0.004}{10.9091}{mi}{86}%
}\hskip.1em\allowbreak{} 的有限维子空间。

证明：本题的“极大无关组”按教材约定，指由有限个向量组成的极大线性无关组。

先证必要性。设\allowbreak{}\hskip.1em \RB{65.681}{8.189}{3.273}{%
\RG{0.000}{-0.002}{10.9091}{mi}{11}\RG{6.982}{-2.697}{7.9701}{rm}{49}\RG{11.825}{-0.002}{10.9091}{rm}{44}\RS{16.669}{-0.002}{10.9091}{\ldots}\RG{27.622}{-0.002}{10.9091}{rm}{44}\RG{32.476}{-0.002}{10.9091}{mi}{11}\RG{39.458}{-2.697}{7.9701}{mi}{114}\RG{47.475}{-0.002}{10.9091}{sy}{50}\RG{58.145}{-0.002}{10.9091}{mi}{83}%
}\hskip.1em\allowbreak{} 构成\allowbreak{}\hskip.1em \RB{7.536}{8.189}{0.742}{%
\RG{0.000}{-0.002}{10.9091}{mi}{83}%
}\hskip.1em\allowbreak{} 的一个极大无关组。任取\allowbreak{}\hskip.1em \RB{28.166}{8.200}{2.618}{%
\RG{0.000}{-0.002}{10.9091}{mi}{12}\RG{9.960}{-0.002}{10.9091}{sy}{50}\RG{20.629}{-0.002}{10.9091}{mi}{83}%
}\hskip.1em\allowbreak{} 。若\allowbreak{}\hskip.1em \RB{6.762}{8.200}{2.618}{%
\RG{0.000}{-0.002}{10.9091}{mi}{12}%
}\hskip.1em\allowbreak{} 不在这些向量中，则由极大性，向量组

\[
\RB{55.686}{8.202}{3.271}{%
\RG{0.000}{-0.000}{10.9091}{mi}{11}\RG{11.825}{-0.000}{10.9091}{rm}{44}\RS{16.669}{-0.000}{10.9091}{\ldots}\RG{27.622}{-0.000}{10.9091}{rm}{44}\RG{32.476}{-0.000}{10.9091}{mi}{11}\RG{44.080}{-0.000}{10.9091}{rm}{44}\RG{48.924}{-0.000}{10.9091}{mi}{12}\RG{6.981}{-2.695}{7.9701}{rm}{49}\RG{39.458}{-2.695}{7.9701}{mi}{114}%
}
\]
线性相关，因而存在不全为零的标量\allowbreak{}\hskip.1em \RB{49.599}{5.327}{3.266}{%
\RG{0.000}{0.005}{10.9091}{mi}{99}\RG{4.724}{-2.690}{7.9701}{rm}{49}\RG{9.567}{0.005}{10.9091}{rm}{44}\RS{14.411}{0.005}{10.9091}{\ldots}\RG{25.364}{0.005}{10.9091}{rm}{44}\RG{30.218}{0.005}{10.9091}{mi}{99}\RG{34.942}{-2.690}{7.9701}{mi}{114}\RG{39.564}{0.005}{10.9091}{rm}{44}\RG{44.408}{0.005}{10.9091}{mi}{99}%
}\hskip.1em\allowbreak{} ，使得

\[
\RB{124.968}{8.199}{3.274}{%
\RG{0.000}{-0.003}{10.9091}{mi}{99}\RG{9.567}{-0.003}{10.9091}{mi}{11}\RG{23.816}{-0.003}{10.9091}{rm}{43}\RS{34.725}{-0.003}{10.9091}{\cdots}\RG{45.613}{-0.003}{10.9091}{rm}{43}\RG{56.522}{-0.003}{10.9091}{mi}{99}\RG{65.876}{-0.003}{10.9091}{mi}{11}\RG{79.904}{-0.003}{10.9091}{rm}{43}\RG{90.813}{-0.003}{10.9091}{mi}{99}\RG{95.809}{-0.003}{10.9091}{mi}{12}\RG{105.409}{-0.003}{10.9091}{rm}{61}\RG{116.919}{-0.003}{10.9091}{rm}{48}\RG{122.373}{-0.003}{10.9091}{rm}{46}\RG{4.723}{-2.698}{7.9701}{rm}{49}\RG{16.548}{-2.698}{7.9701}{rm}{49}\RG{61.254}{-2.698}{7.9701}{mi}{114}\RG{72.858}{-2.698}{7.9701}{mi}{114}%
}
\]
由于\allowbreak{}\hskip.1em \RB{43.868}{5.321}{3.272}{%
\RG{0.000}{-0.001}{10.9091}{mi}{11}\RG{6.982}{-2.696}{7.9701}{rm}{49}\RG{11.825}{-0.001}{10.9091}{rm}{44}\RS{16.669}{-0.001}{10.9091}{\ldots}\RG{27.622}{-0.001}{10.9091}{rm}{44}\RG{32.476}{-0.001}{10.9091}{mi}{11}\RG{39.458}{-2.696}{7.9701}{mi}{114}%
}\hskip.1em\allowbreak{} 线性无关，必有\allowbreak{}\hskip.1em \RB{25.057}{8.463}{3.010}{%
\RG{0.000}{-0.001}{10.9091}{mi}{99}\RS{8.018}{-0.001}{10.9091}{\ne}\RG{19.538}{-0.001}{10.9091}{rm}{48}%
}\hskip.1em\allowbreak{} ，从而

\[
\RB{147.238}{16.494}{13.855}{%
\RG{37.236}{12.626}{7.9701}{mi}{114}\RG{50.188}{8.512}{10.9091}{mi}{99}\RG{54.912}{5.817}{7.9701}{mi}{106}\RG{0.000}{0.004}{10.9091}{mi}{12}\RG{9.589}{0.004}{10.9091}{rm}{61}\RG{21.109}{0.004}{10.9091}{sy}{0}\RG{31.421}{10.369}{10.9092}{ex}{88}\RG{60.218}{0.004}{10.9091}{mi}{11}\RG{74.341}{0.004}{10.9091}{sy}{50}\RG{84.639}{0.004}{10.9091}{mi}{76}\RG{92.068}{0.004}{10.9091}{rm}{40}\RG{96.312}{0.004}{10.9091}{mi}{11}\RG{108.145}{0.004}{10.9091}{rm}{44}\RS{112.989}{0.004}{10.9091}{\ldots}\RG{123.941}{0.004}{10.9091}{rm}{44}\RG{128.796}{0.004}{10.9091}{mi}{11}\RG{140.400}{0.004}{10.9091}{rm}{41}\RG{144.644}{0.004}{10.9091}{rm}{46}\RG{67.200}{-2.690}{7.9701}{mi}{106}\RG{103.302}{-2.690}{7.9701}{rm}{49}\RG{135.778}{-2.690}{7.9701}{mi}{114}\RG{52.107}{-7.480}{10.9091}{mi}{99}\RG{32.349}{-11.797}{7.9701}{mi}{106}\RG{35.961}{-11.797}{7.9701}{rm}{61}\RG{41.902}{-11.797}{7.9701}{rm}{49}\RR{50.188}{2.732}{8.834}{0.436}%
}
\]
当\allowbreak{}\hskip.1em \RB{6.762}{8.207}{2.611}{%
\RG{0.000}{0.005}{10.9091}{mi}{12}%
}\hskip.1em\allowbreak{} 是某个\allowbreak{}\hskip.1em \RB{10.833}{5.327}{4.748}{%
\RG{0.000}{0.005}{10.9091}{mi}{11}\RG{6.981}{-2.690}{7.9701}{mi}{106}%
}\hskip.1em\allowbreak{} 时，该包含关系也成立。因此

\[
\RB{215.396}{9.290}{3.274}{%
\RG{0.000}{-0.003}{10.9091}{mi}{76}\RG{7.429}{-0.003}{10.9091}{rm}{40}\RG{11.673}{-0.003}{10.9091}{mi}{83}\RG{19.004}{-0.003}{10.9091}{rm}{41}\RG{26.280}{-0.003}{10.9091}{rm}{61}\RG{37.800}{-0.003}{10.9091}{mi}{76}\RG{45.229}{-0.003}{10.9091}{rm}{40}\RG{49.473}{-0.003}{10.9091}{mi}{11}\RG{61.305}{-0.003}{10.9091}{rm}{44}\RS{66.148}{-0.003}{10.9091}{\ldots}\RG{77.101}{-0.003}{10.9091}{rm}{44}\RG{81.956}{-0.003}{10.9091}{mi}{11}\RG{93.560}{-0.003}{10.9091}{rm}{41}\RG{97.803}{-0.003}{10.9091}{rm}{44}\RT{124.473}{-0.003}{10.9091}{d}\RT{130.539}{-0.003}{10.9091}{i}\RT{133.571}{-0.003}{10.9091}{m}\RG{144.477}{-0.003}{10.9091}{mi}{76}\RG{151.906}{-0.003}{10.9091}{rm}{40}\RG{156.150}{-0.003}{10.9091}{mi}{83}\RG{163.481}{-0.003}{10.9091}{rm}{41}\RG{170.757}{-0.003}{10.9091}{rm}{61}\RG{182.278}{-0.003}{10.9091}{mi}{114}\RG{190.372}{-0.003}{10.9091}{mi}{60}\RG{201.892}{-0.003}{10.9091}{sy}{49}\RG{212.801}{-0.003}{10.9091}{rm}{46}\RG{56.461}{-2.698}{7.9701}{rm}{49}\RG{88.938}{-2.698}{7.9701}{mi}{114}%
}
\]
再证充分性。设dim\allowbreak{}\hskip.1em \RB{69.075}{8.664}{3.201}{%
\RG{0.000}{0.004}{10.9091}{mi}{76}\RG{7.429}{0.004}{10.9091}{rm}{40}\RG{11.673}{0.004}{10.9091}{mi}{83}\RG{19.004}{0.004}{10.9091}{rm}{41}\RG{26.280}{0.004}{10.9091}{rm}{61}\RG{37.800}{0.004}{10.9091}{mi}{100}\RG{46.768}{0.004}{10.9091}{mi}{60}\RG{58.288}{0.004}{10.9091}{sy}{49}%
}\hskip.1em\allowbreak{} 。由\allowbreak{}\hskip.1em \RB{7.536}{8.195}{0.736}{%
\RG{0.000}{0.004}{10.9091}{mi}{83}%
}\hskip.1em\allowbreak{} 含非零向量，\allowbreak{}\hskip.1em \RB{7.536}{8.195}{0.736}{%
\RG{0.000}{0.004}{10.9091}{mi}{83}%
}\hskip.1em\allowbreak{} 中存在至少含一个向量的线性无关组；其中任一线性无关组所含向量的个数不超过\allowbreak{}\hskip.1em \RB{6.129}{8.066}{0.625}{%
\RG{0.000}{-0.005}{10.9091}{mi}{100}%
}\hskip.1em\allowbreak{} 。

从\allowbreak{}\hskip.1em \RB{7.536}{8.192}{0.739}{%
\RG{0.000}{0.001}{10.9091}{mi}{83}%
}\hskip.1em\allowbreak{} 的所有有限线性无关组中，取所含向量个数最多的一组

\[
\RB{46.674}{5.325}{3.268}{%
\RG{0.000}{0.003}{10.9091}{mi}{11}\RG{11.825}{0.003}{10.9091}{rm}{44}\RS{16.669}{0.003}{10.9091}{\ldots}\RG{27.622}{0.003}{10.9091}{rm}{44}\RG{32.476}{0.003}{10.9091}{mi}{11}\RG{44.079}{0.003}{10.9091}{rm}{46}\RG{6.981}{-2.692}{7.9701}{rm}{49}\RG{39.458}{-2.692}{7.9701}{mi}{114}%
}
\]
若存在\allowbreak{}\hskip.1em \RB{27.435}{8.197}{2.621}{%
\RG{0.000}{-0.005}{10.9091}{mi}{12}\RG{9.589}{-0.005}{10.9091}{sy}{50}\RG{19.898}{-0.005}{10.9091}{mi}{83}%
}\hskip.1em\allowbreak{} ，使\allowbreak{}\hskip.1em \RB{55.686}{8.197}{3.275}{%
\RG{0.000}{-0.005}{10.9091}{mi}{11}\RG{6.982}{-2.699}{7.9701}{rm}{49}\RG{11.826}{-0.005}{10.9091}{rm}{44}\RS{16.670}{-0.005}{10.9091}{\ldots}\RG{27.623}{-0.005}{10.9091}{rm}{44}\RG{32.466}{-0.005}{10.9091}{mi}{11}\RG{39.459}{-2.699}{7.9701}{mi}{114}\RG{44.080}{-0.005}{10.9091}{rm}{44}\RG{48.924}{-0.005}{10.9091}{mi}{12}%
}\hskip.1em\allowbreak{} 仍线性无关，就与\allowbreak{}\hskip.1em \RB{5.256}{5.317}{0.625}{%
\RG{0.000}{-0.005}{10.9091}{mi}{114}%
}\hskip.1em\allowbreak{} 的最大性矛盾。因此该向量组是\allowbreak{}\hskip.1em \RB{7.536}{8.186}{0.745}{%
\RG{0.000}{-0.005}{10.9091}{mi}{83}%
}\hskip.1em\allowbreak{} 的一个极大无关组。

小结：有限极大无关组就是张成空间的一组基；有限维性则保证选取线性无关向量的过程必然终止。

\[
\RB{278.451}{8.393}{2.237}{%
\RT{0.000}{0.002}{7.9701}{题}\RT{8.966}{0.002}{7.9701}{目}\RT{17.933}{0.002}{7.9701}{来}\RT{26.899}{0.002}{7.9701}{源}\RT{35.866}{0.002}{7.9701}{：}\RT{44.832}{0.002}{7.9701}{谢}\RT{53.799}{0.002}{7.9701}{启}\RT{62.765}{0.002}{7.9701}{鸿}\RT{71.732}{0.002}{7.9701}{《}\RT{80.698}{0.002}{7.9701}{复}\RT{89.665}{0.002}{7.9701}{旦}\RT{98.631}{0.002}{7.9701}{大}\RT{107.598}{0.002}{7.9701}{学}\RT{116.564}{0.002}{7.9701}{高}\RT{125.531}{0.002}{7.9701}{等}\RT{134.497}{0.002}{7.9701}{代}\RT{143.464}{0.002}{7.9701}{数}\RT{152.430}{0.002}{7.9701}{习}\RT{161.397}{0.002}{7.9701}{题}\RT{170.363}{0.002}{7.9701}{集}\RT{179.329}{0.002}{7.9701}{（}\RT{188.296}{0.002}{7.9701}{2}\RT{192.779}{0.002}{7.9701}{0}\RT{197.262}{0.002}{7.9701}{2}\RT{201.745}{0.002}{7.9701}{3}\RT{209.214}{0.002}{7.9701}{版}\RT{218.180}{0.002}{7.9701}{）}\RT{222.672}{0.002}{7.9701}{》}\RT{227.156}{0.002}{7.9701}{，}\RT{236.122}{0.002}{7.9701}{第}\RT{248.066}{0.002}{7.9701}{6}\RT{252.549}{0.002}{7.9701}{8}\RT{260.018}{0.002}{7.9701}{页}\RT{268.985}{0.002}{7.9701}{。}%
}
\]
% Original page 3.
\section*{2018A10}
设\allowbreak{}\hskip.1em \RB{22.329}{7.956}{2.600}{%
\RG{0.000}{0.005}{10.9091}{mi}{86}\RG{8.684}{0.005}{10.9091}{rm}{44}\RG{13.538}{0.005}{10.9091}{mi}{85}%
}\hskip.1em\allowbreak{} 分别是数域\allowbreak{}\hskip.1em \RB{6.860}{7.956}{0.500}{%
\RG{0.000}{0.005}{10.9091}{bb}{75}%
}\hskip.1em\allowbreak{} 上的\allowbreak{}\hskip.1em \RB{21.140}{5.327}{2.600}{%
\RG{0.000}{0.005}{10.9091}{mi}{110}\RG{6.546}{0.005}{10.9091}{rm}{44}\RG{11.389}{0.005}{10.9091}{mi}{109}%
}\hskip.1em\allowbreak{} 维线性空间，\allowbreak{}\hskip.1em \RB{73.330}{8.076}{2.873}{%
\RG{0.000}{0.005}{10.9091}{mi}{30}\RG{7.135}{0.005}{10.9091}{rm}{44}\RG{11.978}{0.005}{10.9091}{mi}{32}\RG{24.884}{0.005}{10.9091}{rm}{58}\RG{33.589}{0.005}{10.9091}{mi}{86}\RG{47.957}{0.005}{10.9091}{sy}{33}\RG{64.539}{0.005}{10.9091}{mi}{85}%
}\hskip.1em\allowbreak{} 是两个线性映射。证明：Im\allowbreak{}\hskip.1em \RB{18.217}{7.336}{2.882}{%
\RG{0.000}{-0.004}{10.9091}{mi}{30}\RG{10.156}{-0.004}{10.9091}{sy}{18}%
}\hskip.1em\allowbreak{} Im\allowbreak{}\hskip.1em \RB{7.427}{8.067}{2.740}{%
\RG{0.000}{-0.004}{10.9091}{mi}{32}%
}\hskip.1em\allowbreak{} 的充分必要条件是存在\allowbreak{}\hskip.1em \RB{8.889}{7.947}{0.744}{%
\RG{0.000}{-0.004}{10.9091}{mi}{86}%
}\hskip.1em\allowbreak{} 上的线性变换\allowbreak{}\hskip.1em \RB{5.365}{8.100}{2.740}{%
\RG{0.000}{-0.004}{10.9091}{mi}{24}%
}\hskip.1em\allowbreak{} ，使得

\[
\RB{36.675}{8.102}{2.880}{%
\RG{0.000}{-0.002}{10.9091}{mi}{30}\RG{10.156}{-0.002}{10.9091}{rm}{61}\RG{21.677}{-0.002}{10.9091}{mi}{32}\RG{28.909}{-0.002}{10.9091}{mi}{24}\RG{34.080}{-0.002}{10.9091}{rm}{46}%
}
\]
证明：若存在线性变换\allowbreak{}\hskip.1em \RB{48.817}{8.099}{2.741}{%
\RG{0.000}{-0.005}{10.9091}{mi}{24}\RG{8.193}{-0.005}{10.9091}{rm}{58}\RG{14.258}{-0.005}{10.9091}{mi}{86}\RG{25.986}{-0.005}{10.9091}{sy}{33}\RG{39.928}{-0.005}{10.9091}{mi}{86}%
}\hskip.1em\allowbreak{} ，使得\allowbreak{}\hskip.1em \RB{34.275}{8.099}{2.883}{%
\RG{0.000}{-0.005}{10.9091}{mi}{30}\RG{10.156}{-0.005}{10.9091}{rm}{61}\RG{21.677}{-0.005}{10.9091}{mi}{32}\RG{28.909}{-0.005}{10.9091}{mi}{24}%
}\hskip.1em\allowbreak{} ，则对任意\allowbreak{}\hskip.1em \RB{27.631}{7.946}{0.974}{%
\RG{0.000}{-0.005}{10.9091}{mi}{118}\RG{8.433}{-0.005}{10.9091}{sy}{50}\RG{18.742}{-0.005}{10.9091}{mi}{86}%
}\hskip.1em\allowbreak{} ，

\[
\RB{108.489}{9.290}{3.208}{%
\RG{0.000}{-0.003}{10.9091}{mi}{30}\RG{7.135}{-0.003}{10.9091}{rm}{40}\RG{11.378}{-0.003}{10.9091}{mi}{118}\RG{16.789}{-0.003}{10.9091}{rm}{41}\RG{24.055}{-0.003}{10.9091}{rm}{61}\RG{35.575}{-0.003}{10.9091}{mi}{32}\RG{42.808}{-0.003}{10.9091}{rm}{40}\RG{47.051}{-0.003}{10.9091}{mi}{24}\RG{52.222}{-0.003}{10.9091}{rm}{40}\RG{56.466}{-0.003}{10.9091}{mi}{118}\RG{61.877}{-0.003}{10.9091}{rm}{41}\RG{66.121}{-0.003}{10.9091}{rm}{41}\RG{73.397}{-0.003}{10.9091}{sy}{50}\RT{83.709}{-0.003}{10.9091}{I}\RT{87.647}{-0.003}{10.9091}{m}\RG{98.553}{-0.003}{10.9091}{mi}{32}\RG{105.775}{-0.003}{10.9091}{rm}{44}%
}
\]
故Im\allowbreak{}\hskip.1em \RB{18.217}{7.339}{2.879}{%
\RG{0.000}{-0.001}{10.9091}{mi}{30}\RG{10.156}{-0.001}{10.9091}{sy}{18}%
}\hskip.1em\allowbreak{} Im\allowbreak{}\hskip.1em \RB{7.427}{8.070}{2.737}{%
\RG{0.000}{-0.001}{10.9091}{mi}{32}%
}\hskip.1em\allowbreak{} 。

反之，设Im\allowbreak{}\hskip.1em \RB{18.522}{7.345}{2.873}{%
\RG{0.000}{0.005}{10.9091}{mi}{30}\RG{10.462}{0.005}{10.9091}{sy}{18}%
}\hskip.1em\allowbreak{} Im\allowbreak{}\hskip.1em \RB{7.427}{8.076}{2.731}{%
\RG{0.000}{0.005}{10.9091}{mi}{32}%
}\hskip.1em\allowbreak{} 。取\allowbreak{}\hskip.1em \RB{8.889}{7.956}{0.735}{%
\RG{0.000}{0.005}{10.9091}{mi}{86}%
}\hskip.1em\allowbreak{} 的一组基\allowbreak{}\hskip.1em \RB{41.187}{5.327}{3.266}{%
\RG{0.000}{0.005}{10.9091}{mi}{101}\RG{5.083}{-2.690}{7.9701}{rm}{49}\RG{9.927}{0.005}{10.9091}{rm}{44}\RS{14.771}{0.005}{10.9091}{\ldots}\RG{25.724}{0.005}{10.9091}{rm}{44}\RG{30.578}{0.005}{10.9091}{mi}{101}\RG{35.662}{-2.690}{7.9701}{mi}{110}%
}\hskip.1em\allowbreak{} 。对每个\allowbreak{}\hskip.1em \RB{3.696}{7.716}{0.615}{%
\RG{0.000}{0.005}{10.9091}{mi}{105}%
}\hskip.1em\allowbreak{} ，由\allowbreak{}\hskip.1em \RB{34.523}{8.665}{3.266}{%
\RG{0.000}{0.005}{10.9091}{mi}{30}\RG{7.135}{0.005}{10.9091}{rm}{40}\RG{11.378}{0.005}{10.9091}{mi}{101}\RG{16.462}{-2.690}{7.9701}{mi}{105}\RG{20.045}{0.005}{10.9091}{rm}{41}\RG{27.616}{0.005}{10.9091}{sy}{50}%
}\hskip.1em\allowbreak{} Im\allowbreak{}\hskip.1em \RB{7.427}{8.076}{2.731}{%
\RG{0.000}{0.005}{10.9091}{mi}{32}%
}\hskip.1em\allowbreak{} ，可取\allowbreak{}\hskip.1em \RB{31.702}{7.956}{3.266}{%
\RG{0.000}{0.005}{10.9091}{mi}{118}\RG{5.291}{-2.690}{7.9701}{mi}{105}\RG{12.209}{0.005}{10.9091}{sy}{50}\RG{22.813}{0.005}{10.9091}{mi}{86}%
}\hskip.1em\allowbreak{} ，使得

\[
\RB{137.547}{8.657}{3.274}{%
\RG{0.000}{-0.003}{10.9091}{mi}{32}\RG{7.233}{-0.003}{10.9091}{rm}{40}\RG{11.476}{-0.003}{10.9091}{mi}{118}\RG{20.351}{-0.003}{10.9091}{rm}{41}\RG{27.617}{-0.003}{10.9091}{rm}{61}\RG{39.137}{-0.003}{10.9091}{mi}{30}\RG{46.271}{-0.003}{10.9091}{rm}{40}\RG{50.515}{-0.003}{10.9091}{mi}{101}\RG{59.189}{-0.003}{10.9091}{rm}{41}\RG{63.432}{-0.003}{10.9091}{rm}{44}\RG{90.094}{-0.003}{10.9091}{rm}{49}\RG{98.582}{-0.003}{10.9091}{sy}{20}\RG{110.102}{-0.003}{10.9091}{mi}{105}\RG{116.898}{-0.003}{10.9091}{sy}{20}\RG{128.407}{-0.003}{10.9091}{mi}{110}\RG{134.953}{-0.003}{10.9091}{rm}{46}\RG{16.767}{-2.698}{7.9701}{mi}{105}\RG{55.604}{-2.698}{7.9701}{mi}{105}%
}
\]
\noindent \RB{289.091}{30.992}{14.084}{%
\RT{0.000}{20.892}{10.9091}{定}\RT{10.909}{20.892}{10.9091}{义}\RG{195.254}{12.620}{7.9701}{mi}{110}\RG{255.409}{12.620}{7.9701}{mi}{110}\RG{175.356}{-0.002}{10.9091}{mi}{24}\RG{180.516}{18.599}{10.8765}{ex}{32}\RG{190.073}{10.362}{10.9092}{ex}{88}\RG{207.644}{-0.002}{10.9091}{mi}{120}\RG{217.467}{-0.002}{10.9091}{mi}{101}\RG{226.134}{18.599}{10.8765}{ex}{33}\RG{238.701}{-0.002}{10.9091}{rm}{61}\RG{250.228}{10.362}{10.9092}{ex}{88}\RG{267.799}{-0.002}{10.9091}{mi}{120}\RG{277.622}{-0.002}{10.9091}{mi}{118}\RG{286.496}{-0.002}{10.9091}{rm}{46}\RG{213.884}{-2.697}{7.9701}{mi}{105}\RG{222.551}{-2.697}{7.9701}{mi}{105}\RG{274.039}{-2.697}{7.9701}{mi}{105}\RG{282.913}{-2.697}{7.9701}{mi}{105}\RG{191.264}{-11.803}{7.9701}{mi}{105}\RG{194.349}{-11.803}{7.9701}{rm}{61}\RG{200.290}{-11.803}{7.9701}{rm}{49}\RG{251.419}{-11.803}{7.9701}{mi}{105}\RG{254.504}{-11.803}{7.9701}{rm}{61}\RG{260.445}{-11.803}{7.9701}{rm}{49}%
}

基下的坐标表示唯一，故\allowbreak{}\hskip.1em \RB{5.365}{8.100}{2.740}{%
\RG{0.000}{-0.004}{10.9091}{mi}{24}%
}\hskip.1em\allowbreak{} 定义良好；由定义可知它保持向量的加法与数乘，因此是\allowbreak{}\hskip.1em \RB{8.889}{7.947}{0.744}{%
\RG{0.000}{-0.004}{10.9091}{mi}{86}%
}\hskip.1em\allowbreak{} 上的线性变换。

对任意\allowbreak{}\hskip.1em \RB{87.876}{10.144}{4.827}{%
\RG{0.000}{0.003}{10.9091}{mi}{118}\RG{8.433}{0.003}{10.9091}{rm}{61}\RG{19.953}{8.185}{10.9092}{ex}{80}\RG{31.479}{6.276}{7.9701}{mi}{110}\RG{31.479}{-4.251}{7.9701}{mi}{105}\RG{34.564}{-4.251}{7.9701}{rm}{61}\RG{40.506}{-4.251}{7.9701}{rm}{49}\RG{47.167}{0.003}{10.9091}{mi}{120}\RG{53.407}{-2.692}{7.9701}{mi}{105}\RG{56.990}{0.003}{10.9091}{mi}{101}\RG{62.075}{-2.692}{7.9701}{mi}{105}\RG{68.689}{0.003}{10.9091}{sy}{50}\RG{78.987}{0.003}{10.9091}{mi}{86}%
}\hskip.1em\allowbreak{} ，有

\[
\RB{139.191}{19.538}{83.705}{%
\RG{71.297}{12.622}{7.9701}{mi}{110}\RG{0.000}{-0.000}{10.9091}{rm}{40}\RG{4.244}{-0.000}{10.9091}{mi}{32}\RG{11.466}{-0.000}{10.9091}{mi}{24}\RG{16.637}{-0.000}{10.9091}{rm}{41}\RG{20.880}{-0.000}{10.9091}{rm}{40}\RG{25.124}{-0.000}{10.9091}{mi}{118}\RG{30.535}{-0.000}{10.9091}{rm}{41}\RG{37.811}{-0.000}{10.9091}{rm}{61}\RG{49.331}{-0.000}{10.9091}{mi}{32}\RG{56.564}{18.601}{10.8765}{ex}{32}\RG{66.116}{10.365}{10.9092}{ex}{88}\RG{83.687}{-0.000}{10.9091}{mi}{120}\RG{93.511}{-0.000}{10.9091}{mi}{118}\RG{102.385}{18.602}{10.8765}{ex}{33}\RG{89.927}{-2.694}{7.9701}{mi}{105}\RG{98.802}{-2.694}{7.9701}{mi}{105}\RG{67.307}{-11.801}{7.9701}{mi}{105}\RG{70.392}{-11.801}{7.9701}{rm}{61}\RG{76.333}{-11.801}{7.9701}{rm}{49}\RG{54.518}{-23.239}{7.9701}{mi}{110}\RG{37.819}{-35.861}{10.9091}{rm}{61}\RG{49.337}{-25.498}{10.9092}{ex}{88}\RG{66.909}{-35.861}{10.9091}{mi}{120}\RG{76.733}{-35.861}{10.9091}{mi}{32}\RG{83.955}{-35.861}{10.9091}{rm}{40}\RG{88.198}{-35.861}{10.9091}{mi}{118}\RG{97.083}{-35.861}{10.9091}{rm}{41}\RG{73.149}{-38.556}{7.9701}{mi}{105}\RG{93.500}{-38.556}{7.9701}{mi}{105}\RG{50.528}{-47.662}{7.9701}{mi}{105}\RG{53.614}{-47.662}{7.9701}{rm}{61}\RG{59.555}{-47.662}{7.9701}{rm}{49}\RG{54.518}{-58.705}{7.9701}{mi}{110}\RG{37.819}{-71.328}{10.9091}{rm}{61}\RG{49.337}{-60.964}{10.9092}{ex}{88}\RG{66.909}{-71.328}{10.9091}{mi}{120}\RG{76.732}{-71.328}{10.9091}{mi}{30}\RG{83.866}{-71.328}{10.9091}{rm}{40}\RG{88.110}{-71.328}{10.9091}{mi}{101}\RG{96.778}{-71.328}{10.9091}{rm}{41}\RG{104.043}{-71.328}{10.9091}{rm}{61}\RG{115.563}{-71.328}{10.9091}{mi}{30}\RG{122.698}{-71.328}{10.9091}{rm}{40}\RG{126.942}{-71.328}{10.9091}{mi}{118}\RG{132.353}{-71.328}{10.9091}{rm}{41}\RG{136.596}{-71.328}{10.9091}{rm}{46}\RG{73.149}{-74.023}{7.9701}{mi}{105}\RG{93.194}{-74.023}{7.9701}{mi}{105}\RG{50.528}{-83.129}{7.9701}{mi}{105}\RG{53.614}{-83.129}{7.9701}{rm}{61}\RG{59.555}{-83.129}{7.9701}{rm}{49}%
}
\]
所以\allowbreak{}\hskip.1em \RB{34.275}{8.102}{2.880}{%
\RG{0.000}{-0.002}{10.9091}{mi}{30}\RG{10.156}{-0.002}{10.9091}{rm}{61}\RG{21.677}{-0.002}{10.9091}{mi}{32}\RG{28.909}{-0.002}{10.9091}{mi}{24}%
}\hskip.1em\allowbreak{} 。

小结：在一组基上逐个选取原像，再作线性延拓，即可把像空间的包含关系转化为映射的分解：

\[
\RB{278.451}{8.392}{2.238}{%
\RT{0.000}{0.002}{7.9701}{题}\RT{8.966}{0.002}{7.9701}{目}\RT{17.933}{0.002}{7.9701}{来}\RT{26.899}{0.002}{7.9701}{源}\RT{35.866}{0.002}{7.9701}{：}\RT{44.832}{0.002}{7.9701}{谢}\RT{53.799}{0.002}{7.9701}{启}\RT{62.765}{0.002}{7.9701}{鸿}\RT{71.732}{0.002}{7.9701}{《}\RT{80.698}{0.002}{7.9701}{复}\RT{89.665}{0.002}{7.9701}{旦}\RT{98.631}{0.002}{7.9701}{大}\RT{107.598}{0.002}{7.9701}{学}\RT{116.564}{0.002}{7.9701}{高}\RT{125.531}{0.002}{7.9701}{等}\RT{134.497}{0.002}{7.9701}{代}\RT{143.464}{0.002}{7.9701}{数}\RT{152.430}{0.002}{7.9701}{习}\RT{161.397}{0.002}{7.9701}{题}\RT{170.363}{0.002}{7.9701}{集}\RT{179.329}{0.002}{7.9701}{（}\RT{188.296}{0.002}{7.9701}{2}\RT{192.779}{0.002}{7.9701}{0}\RT{197.262}{0.002}{7.9701}{2}\RT{201.745}{0.002}{7.9701}{3}\RT{209.214}{0.002}{7.9701}{版}\RT{218.180}{0.002}{7.9701}{）}\RT{222.672}{0.002}{7.9701}{》}\RT{227.156}{0.002}{7.9701}{，}\RT{236.122}{0.002}{7.9701}{第}\RT{248.066}{0.002}{7.9701}{6}\RT{252.549}{0.002}{7.9701}{8}\RT{260.018}{0.002}{7.9701}{页}\RT{268.985}{0.002}{7.9701}{。}%
}
\]
% Original page 4.
\section*{2018A11}
(1)请用相抵标准型理论证明：若\allowbreak{}\hskip.1em \RB{8.366}{8.308}{0.503}{%
\RG{0.000}{-0.003}{10.9091}{mi}{65}%
}\hskip.1em\allowbreak{} 为\allowbreak{}\hskip.1em \RB{6.729}{5.319}{0.623}{%
\RG{0.000}{-0.003}{10.9091}{mi}{110}%
}\hskip.1em\allowbreak{} 阶幂等阵，即\allowbreak{}\hskip.1em \RB{35.931}{9.529}{0.503}{%
\RG{0.000}{-0.003}{10.9091}{mi}{65}\RG{8.182}{3.958}{7.9701}{rm}{50}\RG{16.056}{-0.003}{10.9091}{rm}{61}\RG{27.565}{-0.003}{10.9091}{mi}{65}%
}\hskip.1em\allowbreak{} ，则tr\allowbreak{}\hskip.1em \RB{28.068}{8.657}{3.208}{%
\RG{0.000}{-0.003}{10.9091}{rm}{40}\RG{4.244}{-0.003}{10.9091}{mi}{65}\RG{12.426}{-0.003}{10.9091}{rm}{41}\RG{19.691}{-0.003}{10.9091}{rm}{61}%
}\hskip.1em\allowbreak{} r\allowbreak{}\hskip.1em \RB{16.067}{8.657}{3.208}{%
\RG{0.000}{-0.003}{10.9091}{rm}{40}\RG{4.244}{-0.003}{10.9091}{mi}{65}\RG{12.426}{-0.003}{10.9091}{rm}{41}%
}\hskip.1em\allowbreak{} 。

(2)设\allowbreak{}\hskip.1em \RB{7.253}{5.318}{2.882}{%
\RG{0.000}{-0.004}{10.9091}{mi}{30}%
}\hskip.1em\allowbreak{} 是\allowbreak{}\hskip.1em \RB{6.729}{5.318}{0.624}{%
\RG{0.000}{-0.004}{10.9091}{mi}{110}%
}\hskip.1em\allowbreak{} 维线性空间\allowbreak{}\hskip.1em \RB{8.889}{7.947}{0.744}{%
\RG{0.000}{-0.004}{10.9091}{mi}{86}%
}\hskip.1em\allowbreak{} 上的线性变换，满足\allowbreak{}\hskip.1em \RB{41.799}{7.947}{3.352}{%
\RG{0.000}{-0.004}{10.9091}{mi}{30}\RG{7.134}{3.957}{7.9701}{mi}{109}\RG{18.406}{-0.004}{10.9091}{rm}{61}\RG{29.915}{-0.004}{10.9091}{mi}{73}\RG{34.724}{-2.699}{7.9701}{mi}{86}%
}\hskip.1em\allowbreak{} （\allowbreak{}\hskip.1em \RB{29.518}{7.762}{1.802}{%
\RG{0.000}{-0.004}{10.9091}{mi}{109}\RG{12.600}{-0.004}{10.9091}{sy}{21}\RG{24.120}{-0.004}{10.9091}{rm}{50}%
}\hskip.1em\allowbreak{} ），\allowbreak{}\hskip.1em \RB{23.137}{7.947}{0.744}{%
\RG{0.000}{-0.004}{10.9091}{mi}{87}\RG{14.760}{-0.004}{10.9091}{rm}{61}%
}\hskip.1em\allowbreak{} Ker\allowbreak{}\hskip.1em \RB{40.055}{8.656}{3.352}{%
\RG{0.000}{-0.004}{10.9091}{rm}{40}\RG{4.244}{-0.004}{10.9091}{mi}{73}\RG{9.044}{-2.699}{7.9701}{mi}{86}\RG{18.566}{-0.004}{10.9091}{sy}{0}\RG{29.279}{-0.004}{10.9091}{mi}{30}\RG{36.414}{-0.004}{10.9091}{rm}{41}%
}\hskip.1em\allowbreak{} 。证明：

线性变换

\[
\RB{66.477}{18.194}{12.453}{%
\RG{36.322}{12.623}{7.9701}{mi}{109}\RG{44.066}{12.623}{7.9701}{sy}{0}\RG{50.007}{12.623}{7.9701}{rm}{49}\RG{25.791}{7.388}{10.9091}{rm}{49}\RG{63.305}{3.962}{7.9701}{mi}{105}\RG{0.000}{0.001}{10.9091}{mi}{84}\RG{11.007}{0.001}{10.9091}{rm}{61}\RG{37.460}{10.366}{10.9092}{ex}{88}\RG{56.169}{0.001}{10.9091}{mi}{30}\RG{23.730}{-7.483}{10.9091}{mi}{109}\RG{38.651}{-11.800}{7.9701}{mi}{105}\RG{41.736}{-11.800}{7.9701}{rm}{61}\RG{47.678}{-11.800}{7.9701}{rm}{48}\RR{23.730}{2.729}{9.578}{0.436}%
}
\]
的迹等于dim\allowbreak{}\hskip.1em \RB{11.933}{7.954}{0.737}{%
\RG{0.000}{0.003}{10.9091}{mi}{87}%
}\hskip.1em\allowbreak{} 。

证明：(1)记\allowbreak{}\hskip.1em \RB{16.460}{5.324}{0.618}{%
\RG{0.000}{0.002}{10.9091}{mi}{114}\RG{8.084}{0.002}{10.9091}{rm}{61}%
}\hskip.1em\allowbreak{} r\allowbreak{}\hskip.1em \RB{16.067}{8.662}{3.203}{%
\RG{0.000}{0.002}{10.9091}{rm}{40}\RG{4.244}{0.002}{10.9091}{mi}{65}\RG{12.426}{0.002}{10.9091}{rm}{41}%
}\hskip.1em\allowbreak{} 。由相抵标准型理论，存在非异阵\allowbreak{}\hskip.1em \RB{21.947}{8.193}{2.614}{%
\RG{0.000}{0.002}{10.9091}{mi}{80}\RG{8.531}{0.002}{10.9091}{rm}{44}\RG{13.375}{0.002}{10.9091}{mi}{81}%
}\hskip.1em\allowbreak{} ，使得

\[
\RB{150.521}{19.540}{14.080}{%
\RG{108.807}{7.608}{10.9091}{mi}{73}\RG{128.192}{7.608}{10.9091}{mi}{79}\RG{113.607}{4.913}{7.9701}{mi}{114}\RG{0.000}{0.002}{10.9091}{mi}{65}\RG{11.204}{0.002}{10.9091}{rm}{61}\RG{22.724}{0.002}{10.9091}{mi}{80}\RG{31.255}{0.002}{10.9091}{mi}{68}\RG{40.331}{0.002}{10.9091}{mi}{81}\RG{48.960}{0.002}{10.9091}{rm}{44}\RG{75.633}{0.002}{10.9091}{mi}{68}\RG{87.743}{0.002}{10.9091}{rm}{61}\RG{99.252}{18.603}{10.8765}{ex}{32}\RG{136.570}{18.604}{10.8765}{ex}{33}\RG{147.927}{0.003}{10.9091}{rm}{46}\RG{109.329}{-8.651}{10.9091}{mi}{79}\RG{128.191}{-8.651}{10.9091}{mi}{79}%
}
\]
将其代入\allowbreak{}\hskip.1em \RB{35.931}{9.534}{0.500}{%
\RG{0.000}{0.002}{10.9091}{mi}{65}\RG{8.182}{3.963}{7.9701}{rm}{50}\RG{16.056}{0.002}{10.9091}{rm}{61}\RG{27.565}{0.002}{10.9091}{mi}{65}%
}\hskip.1em\allowbreak{} ，再左乘\allowbreak{}\hskip.1em \RB{18.584}{9.534}{0.500}{%
\RG{0.000}{0.002}{10.9091}{mi}{80}\RG{8.531}{3.963}{7.9701}{sy}{0}\RG{14.472}{3.963}{7.9701}{rm}{49}%
}\hskip.1em\allowbreak{} 、右乘\allowbreak{}\hskip.1em \RB{18.683}{9.534}{2.614}{%
\RG{0.000}{0.002}{10.9091}{mi}{81}\RG{8.630}{3.963}{7.9701}{sy}{0}\RG{14.571}{3.963}{7.9701}{rm}{49}%
}\hskip.1em\allowbreak{} ，得到

\[
\RB{70.013}{8.664}{3.201}{%
\RG{0.000}{0.004}{10.9091}{mi}{68}\RG{9.076}{0.004}{10.9091}{rm}{40}\RG{13.320}{0.004}{10.9091}{mi}{81}\RG{21.949}{0.004}{10.9091}{mi}{80}\RG{30.480}{0.004}{10.9091}{rm}{41}\RG{34.724}{0.004}{10.9091}{mi}{68}\RG{46.822}{0.004}{10.9091}{rm}{61}\RG{58.342}{0.004}{10.9091}{mi}{68}\RG{67.419}{0.004}{10.9091}{rm}{46}%
}
\]
因此\allowbreak{}\hskip.1em \RB{17.355}{8.187}{2.620}{%
\RG{0.000}{-0.004}{10.9091}{mi}{81}\RG{8.629}{-0.004}{10.9091}{mi}{80}%
}\hskip.1em\allowbreak{} 左上角的\allowbreak{}\hskip.1em \RB{23.649}{5.852}{0.624}{%
\RG{0.000}{-0.004}{10.9091}{mi}{114}\RG{7.484}{-0.004}{10.9091}{sy}{2}\RG{18.393}{-0.004}{10.9091}{mi}{114}%
}\hskip.1em\allowbreak{} 子块为\allowbreak{}\hskip.1em \RB{9.210}{7.947}{3.275}{%
\RG{0.000}{-0.004}{10.9091}{mi}{73}\RG{4.800}{-2.699}{7.9701}{mi}{114}%
}\hskip.1em\allowbreak{} ，于是

\[
\RB{162.972}{9.291}{3.207}{%
\RT{0.000}{-0.002}{10.9091}{t}\RT{4.244}{-0.002}{10.9091}{r}\RG{8.520}{-0.002}{10.9091}{rm}{40}\RG{12.764}{-0.002}{10.9091}{mi}{65}\RG{20.946}{-0.002}{10.9091}{rm}{41}\RG{28.211}{-0.002}{10.9091}{rm}{61}\RT{39.737}{-0.002}{10.9091}{t}\RT{43.981}{-0.002}{10.9091}{r}\RG{48.257}{-0.002}{10.9091}{rm}{40}\RG{52.501}{-0.002}{10.9091}{mi}{80}\RG{61.032}{-0.002}{10.9091}{mi}{68}\RG{70.109}{-0.002}{10.9091}{mi}{81}\RG{78.738}{-0.002}{10.9091}{rm}{41}\RG{86.003}{-0.002}{10.9091}{rm}{61}\RT{97.530}{-0.002}{10.9091}{t}\RT{101.773}{-0.002}{10.9091}{r}\RG{106.050}{-0.002}{10.9091}{rm}{40}\RG{110.294}{-0.002}{10.9091}{mi}{68}\RG{119.370}{-0.002}{10.9091}{mi}{81}\RG{127.999}{-0.002}{10.9091}{mi}{80}\RG{136.530}{-0.002}{10.9091}{rm}{41}\RG{143.796}{-0.002}{10.9091}{rm}{61}\RG{155.316}{-0.002}{10.9091}{mi}{114}\RG{160.378}{-0.002}{10.9091}{rm}{46}%
}
\]
这里使用了tr\allowbreak{}\hskip.1em \RB{38.093}{8.661}{3.204}{%
\RG{0.000}{0.001}{10.9091}{rm}{40}\RG{4.244}{0.001}{10.9091}{mi}{88}\RG{13.833}{0.001}{10.9091}{mi}{89}\RG{22.451}{0.001}{10.9091}{rm}{41}\RG{29.717}{0.001}{10.9091}{rm}{61}%
}\hskip.1em\allowbreak{} tr\allowbreak{}\hskip.1em \RB{26.093}{8.661}{3.204}{%
\RG{0.000}{0.001}{10.9091}{rm}{40}\RG{4.244}{0.001}{10.9091}{mi}{89}\RG{12.851}{0.001}{10.9091}{mi}{88}\RG{22.451}{0.001}{10.9091}{rm}{41}%
}\hskip.1em\allowbreak{} ，它直接由

\[
\RB{200.567}{10.867}{13.849}{%
\RT{0.000}{0.003}{10.9091}{t}\RT{4.244}{0.003}{10.9091}{r}\RG{8.520}{0.003}{10.9091}{rm}{40}\RG{12.764}{0.003}{10.9091}{mi}{88}\RG{22.342}{0.003}{10.9091}{mi}{89}\RG{30.960}{0.003}{10.9091}{rm}{41}\RG{38.237}{0.003}{10.9091}{rm}{61}\RG{49.757}{10.367}{10.9092}{ex}{88}\RG{67.335}{0.003}{10.9091}{mi}{120}\RG{80.770}{0.003}{10.9091}{mi}{121}\RG{96.341}{0.003}{10.9091}{rm}{61}\RG{107.850}{10.367}{10.9092}{ex}{88}\RG{125.430}{0.003}{10.9091}{mi}{121}\RG{137.971}{0.003}{10.9091}{mi}{120}\RG{154.436}{0.003}{10.9091}{rm}{61}\RT{165.954}{0.003}{10.9091}{t}\RT{170.198}{0.003}{10.9091}{r}\RG{174.474}{0.003}{10.9091}{rm}{40}\RG{178.718}{0.003}{10.9091}{mi}{89}\RG{187.336}{0.003}{10.9091}{mi}{88}\RG{196.926}{0.003}{10.9091}{rm}{41}\RG{73.575}{-2.692}{7.9701}{mi}{105}\RG{76.660}{-2.692}{7.9701}{mi}{106}\RG{86.115}{-2.692}{7.9701}{mi}{106}\RG{89.727}{-2.692}{7.9701}{mi}{105}\RG{130.775}{-2.692}{7.9701}{mi}{106}\RG{134.387}{-2.692}{7.9701}{mi}{105}\RG{144.211}{-2.692}{7.9701}{mi}{105}\RG{147.296}{-2.692}{7.9701}{mi}{106}\RG{52.996}{-11.791}{7.9701}{mi}{105}\RG{56.082}{-11.791}{7.9701}{rm}{44}\RG{58.670}{-11.791}{7.9701}{mi}{106}\RG{111.092}{-11.791}{7.9701}{mi}{106}\RG{114.704}{-11.791}{7.9701}{rm}{44}\RG{117.293}{-11.791}{7.9701}{mi}{105}%
}
\]
得到；\allowbreak{}\hskip.1em \RB{25.122}{7.769}{0.737}{%
\RG{0.000}{0.003}{10.9091}{mi}{114}\RG{8.084}{0.003}{10.9091}{rm}{61}\RG{19.604}{0.003}{10.9091}{rm}{48}%
}\hskip.1em\allowbreak{} 时\allowbreak{}\hskip.1em \RB{31.297}{8.314}{0.737}{%
\RG{0.000}{0.003}{10.9091}{mi}{65}\RG{11.204}{0.003}{10.9091}{rm}{61}\RG{22.724}{0.003}{10.9091}{mi}{79}%
}\hskip.1em\allowbreak{} ，结论同样成立。

(2)由\allowbreak{}\hskip.1em \RB{41.799}{7.950}{3.349}{%
\RG{0.000}{-0.001}{10.9091}{mi}{30}\RG{7.135}{3.959}{7.9701}{mi}{109}\RG{18.407}{-0.001}{10.9091}{rm}{61}\RG{29.916}{-0.001}{10.9091}{mi}{73}\RG{34.724}{-2.696}{7.9701}{mi}{86}%
}\hskip.1em\allowbreak{} ，

\[
\RB{155.767}{18.196}{12.450}{%
\RG{47.461}{12.625}{7.9701}{mi}{109}\RG{100.082}{12.625}{7.9701}{mi}{109}\RG{107.825}{12.625}{7.9701}{sy}{0}\RG{113.766}{12.625}{7.9701}{rm}{49}\RG{32.926}{7.390}{10.9091}{rm}{49}\RG{89.552}{7.390}{10.9091}{rm}{49}\RG{68.163}{3.964}{7.9701}{mi}{105}\RG{127.064}{3.964}{7.9701}{mi}{105}\RG{0.000}{0.004}{10.9091}{mi}{30}\RG{7.135}{0.004}{10.9091}{mi}{84}\RG{18.142}{0.004}{10.9091}{rm}{61}\RG{43.456}{10.368}{10.9092}{ex}{88}\RG{61.028}{0.004}{10.9091}{mi}{30}\RG{74.777}{0.004}{10.9091}{rm}{61}\RG{101.220}{10.368}{10.9092}{ex}{88}\RG{119.930}{0.004}{10.9091}{mi}{30}\RG{133.678}{0.004}{10.9091}{rm}{61}\RG{145.187}{0.004}{10.9091}{mi}{84}\RG{153.173}{0.004}{10.9091}{rm}{46}\RG{30.864}{-7.481}{10.9091}{mi}{109}\RG{87.490}{-7.481}{10.9091}{mi}{109}\RG{44.647}{-11.797}{7.9701}{mi}{105}\RG{47.733}{-11.797}{7.9701}{rm}{61}\RG{53.674}{-11.797}{7.9701}{rm}{49}\RG{102.411}{-11.797}{7.9701}{mi}{105}\RG{105.496}{-11.797}{7.9701}{rm}{61}\RG{111.437}{-11.797}{7.9701}{rm}{48}\RR{30.864}{2.731}{9.578}{0.436}\RR{87.490}{2.731}{9.578}{0.436}%
}
\]
故Im\allowbreak{}\hskip.1em \RB{34.460}{7.949}{1.887}{%
\RG{0.000}{-0.002}{10.9091}{mi}{84}\RG{11.007}{-0.002}{10.9091}{sy}{18}\RG{22.527}{-0.002}{10.9091}{mi}{87}%
}\hskip.1em\allowbreak{} 。反之，若\allowbreak{}\hskip.1em \RB{33.108}{7.949}{0.971}{%
\RG{0.000}{-0.002}{10.9091}{mi}{119}\RG{10.866}{-0.002}{10.9091}{sy}{50}\RG{21.175}{-0.002}{10.9091}{mi}{87}%
}\hskip.1em\allowbreak{} ，则\allowbreak{}\hskip.1em \RB{46.046}{8.658}{3.207}{%
\RG{0.000}{-0.002}{10.9091}{mi}{30}\RG{7.135}{-0.002}{10.9091}{rm}{40}\RG{11.378}{-0.002}{10.9091}{mi}{119}\RG{19.211}{-0.002}{10.9091}{rm}{41}\RG{26.488}{-0.002}{10.9091}{rm}{61}\RG{38.008}{-0.002}{10.9091}{mi}{119}%
}\hskip.1em\allowbreak{} ，从而

\[
\RB{124.530}{18.189}{12.458}{%
\RG{52.652}{12.618}{7.9701}{mi}{109}\RG{60.396}{12.618}{7.9701}{sy}{0}\RG{66.337}{12.618}{7.9701}{rm}{49}\RG{42.122}{7.382}{10.9091}{rm}{49}\RG{79.636}{3.956}{7.9701}{mi}{105}\RG{0.000}{-0.004}{10.9091}{mi}{84}\RG{7.986}{-0.004}{10.9091}{rm}{40}\RG{12.229}{-0.004}{10.9091}{mi}{119}\RG{20.073}{-0.004}{10.9091}{rm}{41}\RG{27.338}{-0.004}{10.9091}{rm}{61}\RG{53.791}{10.360}{10.9092}{ex}{88}\RG{72.501}{-0.004}{10.9091}{mi}{30}\RG{83.219}{-0.004}{10.9091}{rm}{40}\RG{87.462}{-0.004}{10.9091}{mi}{119}\RG{95.295}{-0.004}{10.9091}{rm}{41}\RG{102.572}{-0.004}{10.9091}{rm}{61}\RG{114.092}{-0.004}{10.9091}{mi}{119}\RG{121.935}{-0.004}{10.9091}{rm}{46}\RG{40.060}{-7.488}{10.9091}{mi}{109}\RG{54.981}{-11.805}{7.9701}{mi}{105}\RG{58.067}{-11.805}{7.9701}{rm}{61}\RG{64.008}{-11.805}{7.9701}{rm}{48}\RR{40.060}{2.724}{9.579}{0.436}%
}
\]
于是\allowbreak{}\hskip.1em \RB{22.820}{7.953}{1.883}{%
\RG{0.000}{0.002}{10.9091}{mi}{87}\RG{14.760}{0.002}{10.9091}{sy}{18}%
}\hskip.1em\allowbreak{} Im\allowbreak{}\hskip.1em \RB{8.180}{7.888}{0.500}{%
\RG{0.000}{0.002}{10.9091}{mi}{84}%
}\hskip.1em\allowbreak{} ，且\allowbreak{}\hskip.1em \RB{8.180}{7.888}{0.500}{%
\RG{0.000}{0.002}{10.9091}{mi}{84}%
}\hskip.1em\allowbreak{} 在\allowbreak{}\hskip.1em \RB{11.933}{7.953}{0.738}{%
\RG{0.000}{0.002}{10.9091}{mi}{87}%
}\hskip.1em\allowbreak{} 上为恒等映射。因此Im\allowbreak{}\hskip.1em \RB{34.460}{7.952}{0.739}{%
\RG{0.000}{0.001}{10.9091}{mi}{84}\RG{11.007}{0.001}{10.9091}{rm}{61}\RG{22.527}{0.001}{10.9091}{mi}{87}%
}\hskip.1em\allowbreak{} ，并且对任意\allowbreak{}\hskip.1em \RB{27.631}{7.952}{0.968}{%
\RG{0.000}{0.001}{10.9091}{mi}{118}\RG{8.433}{0.001}{10.9091}{sy}{50}\RG{18.742}{0.001}{10.9091}{mi}{86}%
}\hskip.1em\allowbreak{} ，

\[
\RB{118.656}{9.529}{3.208}{%
\RG{7.985}{3.958}{7.9701}{rm}{50}\RG{0.000}{-0.002}{10.9091}{mi}{84}\RG{12.828}{-0.002}{10.9091}{rm}{40}\RG{17.072}{-0.002}{10.9091}{mi}{118}\RG{22.483}{-0.002}{10.9091}{rm}{41}\RG{29.748}{-0.002}{10.9091}{rm}{61}\RG{41.268}{-0.002}{10.9091}{mi}{84}\RG{49.254}{-0.002}{10.9091}{rm}{40}\RG{53.498}{-0.002}{10.9091}{mi}{84}\RG{61.483}{-0.002}{10.9091}{rm}{40}\RG{65.727}{-0.002}{10.9091}{mi}{118}\RG{71.138}{-0.002}{10.9091}{rm}{41}\RG{75.381}{-0.002}{10.9091}{rm}{41}\RG{82.658}{-0.002}{10.9091}{rm}{61}\RG{94.178}{-0.002}{10.9091}{mi}{84}\RG{102.163}{-0.002}{10.9091}{rm}{40}\RG{106.407}{-0.002}{10.9091}{mi}{118}\RG{111.818}{-0.002}{10.9091}{rm}{41}\RG{116.062}{-0.002}{10.9091}{rm}{46}%
}
\]
由(1)应用于\allowbreak{}\hskip.1em \RB{8.180}{7.885}{0.500}{%
\RG{0.000}{-0.000}{10.9091}{mi}{84}%
}\hskip.1em\allowbreak{} 在任一组基下的矩阵，得到

\[
\RB{166.558}{9.295}{3.204}{%
\RT{0.000}{0.002}{10.9091}{t}\RT{4.244}{0.002}{10.9091}{r}\RG{8.521}{0.002}{10.9091}{rm}{40}\RG{12.765}{0.002}{10.9091}{mi}{84}\RG{20.739}{0.002}{10.9091}{rm}{41}\RG{28.016}{0.002}{10.9091}{rm}{61}\RT{39.542}{0.002}{10.9091}{r}\RG{43.818}{0.002}{10.9091}{rm}{40}\RG{48.062}{0.002}{10.9091}{mi}{84}\RG{56.048}{0.002}{10.9091}{rm}{41}\RG{63.313}{0.002}{10.9091}{rm}{61}\RT{74.840}{0.002}{10.9091}{d}\RT{80.905}{0.002}{10.9091}{i}\RT{83.938}{0.002}{10.9091}{m}\RT{94.836}{0.002}{10.9091}{I}\RT{98.774}{0.002}{10.9091}{m}\RG{109.688}{0.002}{10.9091}{mi}{84}\RG{120.695}{0.002}{10.9091}{rm}{61}\RT{132.222}{0.002}{10.9091}{d}\RT{138.288}{0.002}{10.9091}{i}\RT{141.320}{0.002}{10.9091}{m}\RG{152.225}{0.002}{10.9091}{mi}{87}\RG{163.964}{0.002}{10.9091}{rm}{46}%
}
\]
小结：幂等条件使相抵分解中的关键子块成为单位阵；周期变换的平均则是到不动点空间的投影。

\[
\RB{278.451}{8.392}{2.238}{%
\RT{0.000}{0.002}{7.9701}{题}\RT{8.966}{0.002}{7.9701}{目}\RT{17.933}{0.002}{7.9701}{来}\RT{26.899}{0.002}{7.9701}{源}\RT{35.866}{0.002}{7.9701}{：}\RT{44.832}{0.002}{7.9701}{谢}\RT{53.799}{0.002}{7.9701}{启}\RT{62.765}{0.002}{7.9701}{鸿}\RT{71.732}{0.002}{7.9701}{《}\RT{80.698}{0.002}{7.9701}{复}\RT{89.665}{0.002}{7.9701}{旦}\RT{98.631}{0.002}{7.9701}{大}\RT{107.598}{0.002}{7.9701}{学}\RT{116.564}{0.002}{7.9701}{高}\RT{125.531}{0.002}{7.9701}{等}\RT{134.497}{0.002}{7.9701}{代}\RT{143.464}{0.002}{7.9701}{数}\RT{152.430}{0.002}{7.9701}{习}\RT{161.397}{0.002}{7.9701}{题}\RT{170.363}{0.002}{7.9701}{集}\RT{179.329}{0.002}{7.9701}{（}\RT{188.296}{0.002}{7.9701}{2}\RT{192.779}{0.002}{7.9701}{0}\RT{197.262}{0.002}{7.9701}{2}\RT{201.745}{0.002}{7.9701}{3}\RT{209.214}{0.002}{7.9701}{版}\RT{218.180}{0.002}{7.9701}{）}\RT{222.672}{0.002}{7.9701}{》}\RT{227.156}{0.002}{7.9701}{，}\RT{236.122}{0.002}{7.9701}{第}\RT{248.066}{0.002}{7.9701}{6}\RT{252.549}{0.002}{7.9701}{8}\RT{260.018}{0.002}{7.9701}{页}\RT{268.985}{0.002}{7.9701}{。}%
}
\]
% Original page 5.
\section*{2018A12}
设循环矩阵

\[
\RB{160.719}{44.605}{33.169}{%
\RG{37.942}{35.214}{10.9091}{mi}{97}\RG{64.709}{35.214}{10.9091}{mi}{97}\RG{85.809}{35.214}{10.9091}{mi}{97}\RS{106.385}{35.214}{10.9091}{\cdots}\RG{129.949}{35.214}{10.9091}{mi}{97}\RG{43.713}{32.520}{7.9701}{rm}{49}\RG{70.480}{32.520}{7.9701}{rm}{50}\RG{91.580}{32.520}{7.9701}{rm}{51}\RG{135.728}{32.520}{7.9701}{mi}{110}\RG{22.730}{43.754}{9.0156}{ex}{48}\RG{146.760}{43.754}{9.0156}{ex}{49}\RG{22.730}{28.088}{8.7625}{ex}{66}\RG{146.760}{28.088}{8.7625}{ex}{67}\RG{37.419}{18.956}{10.9091}{mi}{97}\RG{64.708}{18.956}{10.9091}{mi}{97}\RG{85.809}{18.956}{10.9091}{mi}{97}\RS{106.385}{18.956}{10.9091}{\cdots}\RG{124.810}{18.956}{10.9091}{mi}{97}\RG{22.730}{23.210}{8.7625}{ex}{66}\RG{146.760}{23.210}{8.7625}{ex}{67}\RG{43.190}{16.261}{7.9701}{mi}{110}\RG{70.480}{16.261}{7.9701}{rm}{49}\RG{91.580}{16.261}{7.9701}{rm}{50}\RG{130.584}{16.261}{7.9701}{mi}{110}\RG{135.975}{16.261}{7.9701}{sy}{0}\RG{141.917}{16.261}{7.9701}{rm}{49}\RG{22.730}{18.332}{8.7625}{ex}{66}\RG{146.760}{18.332}{8.7625}{ex}{67}\RG{22.730}{13.455}{8.7625}{ex}{66}\RG{146.760}{13.455}{8.7625}{ex}{67}\RG{0.000}{3.221}{10.9091}{mi}{65}\RG{11.204}{3.221}{10.9091}{rm}{61}\RG{22.730}{8.577}{8.7625}{ex}{66}\RG{146.760}{8.577}{8.7625}{ex}{67}\RG{158.124}{3.221}{10.9091}{rm}{46}\RG{32.276}{2.697}{10.9091}{mi}{97}\RG{64.185}{2.697}{10.9091}{mi}{97}\RG{85.809}{2.697}{10.9091}{mi}{97}\RS{106.385}{2.697}{10.9091}{\cdots}\RG{124.810}{2.697}{10.9091}{mi}{97}\RG{38.047}{0.002}{7.9701}{mi}{110}\RG{43.439}{0.002}{7.9701}{sy}{0}\RG{49.380}{0.002}{7.9701}{rm}{49}\RG{69.957}{0.002}{7.9701}{mi}{110}\RG{91.580}{0.002}{7.9701}{rm}{49}\RG{130.584}{0.002}{7.9701}{mi}{110}\RG{135.975}{0.002}{7.9701}{sy}{0}\RG{141.917}{0.002}{7.9701}{rm}{50}\RG{22.730}{3.699}{8.7625}{ex}{66}\RG{146.760}{3.699}{8.7625}{ex}{67}\RG{22.730}{-1.179}{8.7625}{ex}{66}\RG{146.760}{-1.179}{8.7625}{ex}{67}\RG{22.730}{-6.056}{8.7625}{ex}{66}\RG{146.760}{-6.056}{8.7625}{ex}{67}\RS{42.060}{-13.562}{10.9091}{\vdots}\RS{68.821}{-13.562}{10.9091}{\vdots}\RS{89.919}{-13.562}{10.9091}{\vdots}\RS{134.598}{-13.562}{10.9091}{\vdots}\RG{22.730}{-10.934}{8.7625}{ex}{66}\RG{146.760}{-10.934}{8.7625}{ex}{67}\RG{37.942}{-29.821}{10.9091}{mi}{97}\RG{64.708}{-29.821}{10.9091}{mi}{97}\RG{85.809}{-29.821}{10.9091}{mi}{97}\RS{106.385}{-29.821}{10.9091}{\cdots}\RG{130.472}{-29.821}{10.9091}{mi}{97}\RG{22.730}{-16.350}{9.0156}{ex}{64}\RG{146.760}{-16.351}{9.0156}{ex}{65}\RG{43.713}{-32.516}{7.9701}{rm}{50}\RG{70.479}{-32.516}{7.9701}{rm}{51}\RG{91.580}{-32.516}{7.9701}{rm}{52}\RG{136.251}{-32.516}{7.9701}{rm}{49}%
}
\]
证明：伴随阵\allowbreak{}\hskip.1em \RB{12.020}{8.308}{0.503}{%
\RG{0.000}{-0.003}{10.9091}{mi}{65}\RG{8.182}{3.957}{7.9701}{sy}{3}%
}\hskip.1em\allowbreak{} 也是循环矩阵。

证明：令\allowbreak{}\hskip.1em \RB{43.648}{8.663}{4.750}{%
\RG{0.000}{0.003}{10.9091}{mi}{80}\RG{11.553}{0.003}{10.9091}{rm}{61}\RG{23.073}{0.003}{10.9091}{rm}{40}\RG{27.317}{0.003}{10.9091}{mi}{112}\RG{32.811}{-2.692}{7.9701}{mi}{105}\RG{35.896}{-2.692}{7.9701}{mi}{106}\RG{40.006}{0.003}{10.9091}{rm}{41}%
}\hskip.1em\allowbreak{} 为循环移位矩阵，其中\allowbreak{}\hskip.1em \RB{44.641}{7.769}{4.666}{%
\RG{0.000}{0.003}{10.9091}{mi}{112}\RG{5.487}{-2.692}{7.9701}{mi}{105}\RG{8.572}{-2.692}{7.9701}{rm}{44}\RG{11.161}{-2.692}{7.9701}{mi}{105}\RG{14.246}{-2.692}{7.9701}{rm}{43}\RG{20.187}{-2.692}{7.9701}{rm}{49}\RG{28.061}{0.003}{10.9091}{rm}{61}\RG{39.570}{0.003}{10.9091}{rm}{49}%
}\hskip.1em\allowbreak{} （\allowbreak{}\hskip.1em \RB{45.042}{7.769}{1.795}{%
\RG{0.000}{0.003}{10.9091}{rm}{49}\RG{8.476}{0.003}{10.9091}{sy}{20}\RG{19.997}{0.003}{10.9091}{mi}{105}\RG{26.793}{0.003}{10.9091}{mi}{60}\RG{38.313}{0.003}{10.9091}{mi}{110}%
}\hskip.1em\allowbreak{} ）、\allowbreak{}\hskip.1em \RB{35.332}{7.769}{3.268}{%
\RG{0.000}{0.003}{10.9091}{mi}{112}\RG{5.487}{-2.692}{7.9701}{mi}{110}\RG{10.878}{-2.692}{7.9701}{rm}{49}\RG{18.752}{0.003}{10.9091}{rm}{61}\RG{30.261}{0.003}{10.9091}{rm}{49}%
}\hskip.1em\allowbreak{} ，其余元素为零。于是

\[
\RB{145.773}{9.526}{3.276}{%
\RG{126.882}{3.955}{7.9701}{mi}{110}\RG{132.273}{3.955}{7.9701}{sy}{0}\RG{138.215}{3.955}{7.9701}{rm}{49}\RG{0.000}{-0.005}{10.9091}{mi}{65}\RG{11.204}{-0.005}{10.9091}{rm}{61}\RG{22.724}{-0.005}{10.9091}{mi}{97}\RG{33.344}{-0.005}{10.9091}{mi}{73}\RG{41.493}{-0.005}{10.9091}{rm}{43}\RG{52.403}{-0.005}{10.9091}{mi}{97}\RG{63.023}{-0.005}{10.9091}{mi}{80}\RG{73.975}{-0.005}{10.9091}{rm}{43}\RS{84.885}{-0.005}{10.9091}{\cdots}\RG{95.772}{-0.005}{10.9091}{rm}{43}\RG{106.681}{-0.005}{10.9091}{mi}{97}\RG{118.352}{-0.005}{10.9091}{mi}{80}\RG{143.058}{-0.005}{10.9091}{rm}{44}\RG{28.501}{-2.700}{7.9701}{rm}{49}\RG{58.178}{-2.700}{7.9701}{rm}{50}\RG{112.462}{-2.700}{7.9701}{mi}{110}%
}
\]
所以\allowbreak{}\hskip.1em \RB{48.151}{8.308}{0.503}{%
\RG{0.000}{-0.003}{10.9091}{mi}{65}\RG{8.182}{-0.003}{10.9091}{mi}{80}\RG{19.735}{-0.003}{10.9091}{rm}{61}\RG{31.255}{-0.003}{10.9091}{mi}{80}\RG{39.786}{-0.003}{10.9091}{mi}{65}%
}\hskip.1em\allowbreak{} 。

对任意标量\allowbreak{}\hskip.1em \RB{4.100}{7.332}{0.617}{%
\RG{0.000}{0.003}{10.9091}{mi}{116}%
}\hskip.1em\allowbreak{} ，记\allowbreak{}\hskip.1em \RB{57.903}{8.314}{3.268}{%
\RG{0.000}{0.003}{10.9091}{mi}{65}\RG{8.182}{-2.692}{7.9701}{mi}{116}\RG{15.010}{0.003}{10.9091}{rm}{61}\RG{26.519}{0.003}{10.9091}{mi}{65}\RG{37.134}{0.003}{10.9091}{rm}{43}\RG{48.043}{0.003}{10.9091}{mi}{116}\RG{51.981}{0.003}{10.9091}{mi}{73}%
}\hskip.1em\allowbreak{} ，则\allowbreak{}\hskip.1em \RB{55.173}{8.314}{3.268}{%
\RG{0.000}{0.003}{10.9091}{mi}{65}\RG{8.182}{-2.692}{7.9701}{mi}{116}\RG{11.979}{0.003}{10.9091}{mi}{80}\RG{23.532}{0.003}{10.9091}{rm}{61}\RG{35.052}{0.003}{10.9091}{mi}{80}\RG{43.583}{0.003}{10.9091}{mi}{65}\RG{51.771}{-2.692}{7.9701}{mi}{116}%
}\hskip.1em\allowbreak{} 。当det\allowbreak{}\hskip.1em \RB{32.037}{8.467}{3.268}{%
\RG{0.000}{0.003}{10.9091}{mi}{65}\RG{8.181}{-2.692}{7.9701}{mi}{116}\RS{15.009}{0.003}{10.9091}{\ne}\RG{26.518}{0.003}{10.9091}{rm}{48}%
}\hskip.1em\allowbreak{} 时，\allowbreak{}\hskip.1em \RB{18.235}{9.535}{3.268}{%
\RG{0.000}{0.003}{10.9091}{mi}{65}\RG{8.182}{3.964}{7.9701}{sy}{0}\RG{8.182}{-2.692}{7.9701}{mi}{116}\RG{14.123}{3.964}{7.9701}{rm}{49}%
}\hskip.1em\allowbreak{} 也与\allowbreak{}\hskip.1em \RB{8.726}{7.954}{0.500}{%
\RG{0.000}{0.003}{10.9091}{mi}{80}%
}\hskip.1em\allowbreak{} 交换，故

\[
\RB{214.531}{9.537}{3.266}{%
\RG{8.182}{3.966}{7.9701}{sy}{3}\RG{79.379}{3.966}{7.9701}{sy}{0}\RG{85.320}{3.966}{7.9701}{rm}{49}\RG{165.574}{3.966}{7.9701}{sy}{0}\RG{171.516}{3.966}{7.9701}{rm}{49}\RG{207.621}{3.966}{7.9701}{sy}{3}\RG{0.000}{0.005}{10.9091}{mi}{65}\RG{12.498}{0.005}{10.9091}{mi}{80}\RG{24.051}{0.005}{10.9091}{rm}{61}\RT{35.577}{0.005}{10.9091}{d}\RT{41.643}{0.005}{10.9091}{e}\RT{46.486}{0.005}{10.9091}{t}\RG{50.730}{0.005}{10.9091}{rm}{40}\RG{54.974}{0.005}{10.9091}{mi}{65}\RG{66.954}{0.005}{10.9091}{rm}{41}\RG{71.197}{0.005}{10.9091}{mi}{65}\RG{90.163}{0.005}{10.9091}{mi}{80}\RG{101.716}{0.005}{10.9091}{rm}{61}\RT{113.243}{0.005}{10.9091}{d}\RT{119.308}{0.005}{10.9091}{e}\RT{124.152}{0.005}{10.9091}{t}\RG{128.395}{0.005}{10.9091}{rm}{40}\RG{132.639}{0.005}{10.9091}{mi}{65}\RG{144.618}{0.005}{10.9091}{rm}{41}\RG{148.862}{0.005}{10.9091}{mi}{80}\RG{157.382}{0.005}{10.9091}{mi}{65}\RG{179.389}{0.005}{10.9091}{rm}{61}\RG{190.899}{0.005}{10.9091}{mi}{80}\RG{199.430}{0.005}{10.9091}{mi}{65}\RG{211.937}{0.005}{10.9091}{rm}{46}\RG{8.182}{-2.690}{7.9701}{mi}{116}\RG{63.156}{-2.690}{7.9701}{mi}{116}\RG{79.379}{-2.690}{7.9701}{mi}{116}\RG{140.821}{-2.690}{7.9701}{mi}{116}\RG{165.574}{-2.690}{7.9701}{mi}{116}\RG{207.621}{-2.690}{7.9701}{mi}{116}%
}
\]
多项式det\allowbreak{}\hskip.1em \RB{38.508}{8.657}{3.208}{%
\RG{0.000}{-0.003}{10.9091}{rm}{40}\RG{4.244}{-0.003}{10.9091}{mi}{65}\RG{14.564}{-0.003}{10.9091}{rm}{43}\RG{25.200}{-0.003}{10.9091}{mi}{116}\RG{29.138}{-0.003}{10.9091}{mi}{73}\RG{34.866}{-0.003}{10.9091}{rm}{41}%
}\hskip.1em\allowbreak{} 的首项为\allowbreak{}\hskip.1em \RB{9.463}{7.826}{0.623}{%
\RG{0.000}{-0.003}{10.9091}{mi}{116}\RG{3.938}{3.958}{7.9701}{mi}{110}%
}\hskip.1em\allowbreak{} ，因而仅有有限个根。上式遂对无限多个\allowbreak{}\hskip.1em \RB{4.100}{7.326}{0.623}{%
\RG{0.000}{-0.003}{10.9091}{mi}{116}%
}\hskip.1em\allowbreak{} 成立。又因伴随阵的每个元素都是余子式，\allowbreak{}\hskip.1em \RB{55.017}{8.309}{3.273}{%
\RG{0.000}{-0.002}{10.9091}{mi}{65}\RG{8.182}{3.959}{7.9701}{sy}{3}\RG{8.182}{-2.697}{7.9701}{mi}{116}\RG{12.498}{-0.002}{10.9091}{mi}{80}\RG{23.495}{-0.002}{10.9091}{sy}{0}\RG{34.458}{-0.002}{10.9091}{mi}{80}\RG{42.989}{-0.002}{10.9091}{mi}{65}\RG{51.178}{3.959}{7.9701}{sy}{3}\RG{51.178}{-2.697}{7.9701}{mi}{116}%
}\hskip.1em\allowbreak{} 的每个元素均为\allowbreak{}\hskip.1em \RB{4.100}{7.327}{0.622}{%
\RG{0.000}{-0.002}{10.9091}{mi}{116}%
}\hskip.1em\allowbreak{} 的多项式；这些多项式有无限多个根，故恒为零。取\allowbreak{}\hskip.1em \RB{23.998}{7.765}{0.741}{%
\RG{0.000}{-0.001}{10.9091}{mi}{116}\RG{6.960}{-0.001}{10.9091}{rm}{61}\RG{18.480}{-0.001}{10.9091}{rm}{48}%
}\hskip.1em\allowbreak{} ，得到

\[
\RB{59.202}{8.311}{0.500}{%
\RG{8.182}{3.961}{7.9701}{sy}{3}\RG{52.290}{3.961}{7.9701}{sy}{3}\RG{0.000}{-0.000}{10.9091}{mi}{65}\RG{12.498}{-0.000}{10.9091}{mi}{80}\RG{24.051}{-0.000}{10.9091}{rm}{61}\RG{35.571}{-0.000}{10.9091}{mi}{80}\RG{44.102}{-0.000}{10.9091}{mi}{65}\RG{56.607}{-0.000}{10.9091}{rm}{46}%
}
\]
记\allowbreak{}\hskip.1em \RB{46.807}{8.661}{4.752}{%
\RG{0.000}{0.001}{10.9091}{mi}{65}\RG{8.182}{3.961}{7.9701}{sy}{3}\RG{15.529}{0.001}{10.9091}{rm}{61}\RG{27.038}{0.001}{10.9091}{rm}{40}\RG{31.282}{0.001}{10.9091}{mi}{98}\RG{35.970}{-2.694}{7.9701}{mi}{105}\RG{39.055}{-2.694}{7.9701}{mi}{106}\RG{43.165}{0.001}{10.9091}{rm}{41}%
}\hskip.1em\allowbreak{} ，并将行、列下标按模\allowbreak{}\hskip.1em \RB{6.729}{5.323}{0.619}{%
\RG{0.000}{0.001}{10.9091}{mi}{110}%
}\hskip.1em\allowbreak{} 循环理解：比较上式第\allowbreak{}\hskip.1em \RB{21.129}{8.661}{3.204}{%
\RG{0.000}{0.001}{10.9091}{rm}{40}\RG{4.244}{0.001}{10.9091}{mi}{105}\RG{8.007}{0.001}{10.9091}{rm}{44}\RG{12.851}{0.001}{10.9091}{mi}{106}\RG{17.487}{0.001}{10.9091}{rm}{41}%
}\hskip.1em\allowbreak{} 个元素，得

\[
\RB{147.655}{8.074}{4.750}{%
\RG{0.000}{0.003}{10.9091}{mi}{98}\RG{27.780}{0.003}{10.9091}{rm}{61}\RG{39.289}{0.003}{10.9091}{mi}{98}\RG{64.049}{0.003}{10.9091}{rm}{44}\RG{90.711}{0.003}{10.9091}{rm}{49}\RG{99.198}{0.003}{10.9091}{sy}{20}\RG{110.718}{0.003}{10.9091}{mi}{105}\RG{114.482}{0.003}{10.9091}{rm}{44}\RG{119.325}{0.003}{10.9091}{mi}{106}\RG{126.995}{0.003}{10.9091}{sy}{20}\RG{138.515}{0.003}{10.9091}{mi}{110}\RG{145.060}{0.003}{10.9091}{rm}{46}\RG{4.680}{-2.692}{7.9701}{mi}{105}\RG{7.765}{-2.692}{7.9701}{rm}{44}\RG{10.354}{-2.692}{7.9701}{mi}{106}\RG{13.966}{-2.692}{7.9701}{sy}{0}\RG{19.907}{-2.692}{7.9701}{rm}{49}\RG{43.978}{-2.692}{7.9701}{mi}{105}\RG{47.063}{-2.692}{7.9701}{rm}{43}\RG{53.005}{-2.692}{7.9701}{rm}{49}\RG{57.350}{-2.692}{7.9701}{rm}{44}\RG{59.939}{-2.692}{7.9701}{mi}{106}%
}
\]
因此\allowbreak{}\hskip.1em \RB{12.020}{8.306}{0.505}{%
\RG{0.000}{-0.005}{10.9091}{mi}{65}\RG{8.182}{3.955}{7.9701}{sy}{3}%
}\hskip.1em\allowbreak{} 的每一行均由上一行向右循环移动一位得到，即\allowbreak{}\hskip.1em \RB{12.020}{8.306}{0.505}{%
\RG{0.000}{-0.005}{10.9091}{mi}{65}\RG{8.182}{3.955}{7.9701}{sy}{3}%
}\hskip.1em\allowbreak{} 是循环矩阵。

小结：循环矩阵可由与循环移位矩阵交换来刻画；引入\allowbreak{}\hskip.1em \RB{31.602}{8.313}{1.403}{%
\RG{0.000}{0.002}{10.9091}{mi}{65}\RG{10.713}{0.002}{10.9091}{rm}{43}\RG{21.742}{0.002}{10.9091}{mi}{116}\RG{25.680}{0.002}{10.9091}{mi}{73}%
}\hskip.1em\allowbreak{} ，可把可逆情形的结论延伸到奇异情形。

\[
\RB{295.397}{8.393}{2.237}{%
\RT{0.000}{0.002}{7.9701}{题}\RT{8.966}{0.002}{7.9701}{目}\RT{17.933}{0.002}{7.9701}{来}\RT{26.899}{0.002}{7.9701}{源}\RT{35.866}{0.002}{7.9701}{：}\RT{44.832}{0.002}{7.9701}{谢}\RT{53.799}{0.002}{7.9701}{启}\RT{62.765}{0.002}{7.9701}{鸿}\RT{71.732}{0.002}{7.9701}{《}\RT{80.698}{0.002}{7.9701}{复}\RT{89.665}{0.002}{7.9701}{旦}\RT{98.631}{0.002}{7.9701}{大}\RT{107.598}{0.002}{7.9701}{学}\RT{116.564}{0.002}{7.9701}{高}\RT{125.531}{0.002}{7.9701}{等}\RT{134.497}{0.002}{7.9701}{代}\RT{143.464}{0.002}{7.9701}{数}\RT{152.430}{0.002}{7.9701}{习}\RT{161.397}{0.002}{7.9701}{题}\RT{170.363}{0.002}{7.9701}{集}\RT{179.329}{0.002}{7.9701}{（}\RT{188.296}{0.002}{7.9701}{2}\RT{192.779}{0.002}{7.9701}{0}\RT{197.262}{0.002}{7.9701}{2}\RT{201.745}{0.002}{7.9701}{3}\RT{209.214}{0.002}{7.9701}{版}\RT{218.180}{0.002}{7.9701}{）}\RT{222.672}{0.002}{7.9701}{》}\RT{227.156}{0.002}{7.9701}{，}\RT{236.122}{0.002}{7.9701}{第}\RT{248.066}{0.002}{7.9701}{6}\RT{252.549}{0.002}{7.9701}{8}\RT{257.036}{0.002}{7.9701}{—}\RT{265.012}{0.002}{7.9701}{6}\RT{269.495}{0.002}{7.9701}{9}\RT{276.964}{0.002}{7.9701}{页}\RT{285.931}{0.002}{7.9701}{。}%
}
\]
% Original page 6.
\section*{2018A13}
设复系数多项式\allowbreak{}\hskip.1em \RB{45.500}{8.665}{3.200}{%
\RG{0.000}{0.005}{10.9091}{mi}{102}\RG{6.327}{0.005}{10.9091}{rm}{40}\RG{10.571}{0.005}{10.9091}{mi}{120}\RG{16.811}{0.005}{10.9091}{rm}{41}\RG{21.055}{0.005}{10.9091}{rm}{44}\RG{25.898}{0.005}{10.9091}{mi}{103}\RG{31.375}{0.005}{10.9091}{rm}{40}\RG{35.619}{0.005}{10.9091}{mi}{120}\RG{41.859}{0.005}{10.9091}{rm}{41}%
}\hskip.1em\allowbreak{} 互素。证明：\allowbreak{}\hskip.1em \RB{63.794}{9.536}{3.200}{%
\RG{0.000}{0.005}{10.9091}{mi}{102}\RG{6.316}{0.005}{10.9091}{rm}{40}\RG{10.560}{0.005}{10.9091}{mi}{120}\RG{16.800}{0.005}{10.9091}{rm}{41}\RG{21.054}{3.965}{7.9701}{rm}{50}\RG{28.322}{0.005}{10.9091}{rm}{43}\RG{39.231}{0.005}{10.9091}{mi}{103}\RG{44.708}{0.005}{10.9091}{rm}{40}\RG{48.952}{0.005}{10.9091}{mi}{120}\RG{55.192}{0.005}{10.9091}{rm}{41}\RG{59.438}{3.965}{7.9701}{rm}{50}%
}\hskip.1em\allowbreak{} 的重根必为\allowbreak{}\hskip.1em \RB{71.006}{9.536}{3.200}{%
\RG{0.000}{0.005}{10.9091}{mi}{102}\RG{6.327}{3.965}{7.9701}{sy}{48}\RG{9.933}{0.005}{10.9091}{rm}{40}\RG{14.177}{0.005}{10.9091}{mi}{120}\RG{20.417}{0.005}{10.9091}{rm}{41}\RG{24.660}{3.965}{7.9701}{rm}{50}\RG{31.928}{0.005}{10.9091}{rm}{43}\RG{42.837}{0.005}{10.9091}{mi}{103}\RG{48.316}{3.965}{7.9701}{sy}{48}\RG{51.922}{0.005}{10.9091}{rm}{40}\RG{56.166}{0.005}{10.9091}{mi}{120}\RG{62.406}{0.005}{10.9091}{rm}{41}\RG{66.650}{3.965}{7.9701}{rm}{50}%
}\hskip.1em\allowbreak{} 的根。

证明：记

\[
\RB{104.628}{9.533}{3.203}{%
\RG{58.807}{3.962}{7.9701}{rm}{50}\RG{97.191}{3.962}{7.9701}{rm}{50}\RG{0.000}{0.002}{10.9091}{mi}{70}\RG{8.476}{0.002}{10.9091}{rm}{40}\RG{12.720}{0.002}{10.9091}{mi}{120}\RG{18.960}{0.002}{10.9091}{rm}{41}\RG{26.226}{0.002}{10.9091}{rm}{61}\RG{37.746}{0.002}{10.9091}{mi}{102}\RG{44.073}{0.002}{10.9091}{rm}{40}\RG{48.317}{0.002}{10.9091}{mi}{120}\RG{54.557}{0.002}{10.9091}{rm}{41}\RG{66.075}{0.002}{10.9091}{rm}{43}\RG{76.984}{0.002}{10.9091}{mi}{103}\RG{82.460}{0.002}{10.9091}{rm}{40}\RG{86.704}{0.002}{10.9091}{mi}{120}\RG{92.944}{0.002}{10.9091}{rm}{41}\RG{102.034}{0.002}{10.9091}{rm}{46}%
}
\]
设\allowbreak{}\hskip.1em \RB{5.933}{5.319}{0.623}{%
\RG{0.000}{-0.003}{10.9091}{mi}{97}%
}\hskip.1em\allowbreak{} 是\allowbreak{}\hskip.1em \RB{8.669}{7.915}{0.503}{%
\RG{0.000}{-0.003}{10.9091}{mi}{70}%
}\hskip.1em\allowbreak{} 的重根，则

\[
\RB{123.898}{29.781}{3.200}{%
\RG{58.431}{24.210}{7.9701}{rm}{50}\RG{96.346}{24.210}{7.9701}{rm}{50}\RG{37.846}{20.249}{10.9091}{mi}{102}\RG{44.163}{20.249}{10.9091}{rm}{40}\RG{48.406}{20.249}{10.9091}{mi}{97}\RG{54.177}{20.249}{10.9091}{rm}{41}\RG{65.700}{20.249}{10.9091}{rm}{43}\RG{76.609}{20.249}{10.9091}{mi}{103}\RG{82.085}{20.249}{10.9091}{rm}{40}\RG{86.329}{20.249}{10.9091}{mi}{97}\RG{92.100}{20.249}{10.9091}{rm}{41}\RG{104.220}{20.249}{10.9091}{rm}{61}\RG{115.729}{20.249}{10.9091}{rm}{48}\RG{121.184}{20.249}{10.9091}{rm}{44}\RG{26.913}{3.965}{7.9701}{sy}{48}\RG{83.325}{3.965}{7.9701}{sy}{48}\RG{0.000}{0.005}{10.9091}{mi}{102}\RG{6.316}{0.005}{10.9091}{rm}{40}\RG{10.560}{0.005}{10.9091}{mi}{97}\RG{16.331}{0.005}{10.9091}{rm}{41}\RG{20.575}{0.005}{10.9091}{mi}{102}\RG{30.519}{0.005}{10.9091}{rm}{40}\RG{34.763}{0.005}{10.9091}{mi}{97}\RG{40.534}{0.005}{10.9091}{rm}{41}\RG{47.199}{0.005}{10.9091}{rm}{43}\RG{58.109}{0.005}{10.9091}{mi}{103}\RG{63.585}{0.005}{10.9091}{rm}{40}\RG{67.829}{0.005}{10.9091}{mi}{97}\RG{73.600}{0.005}{10.9091}{rm}{41}\RG{77.843}{0.005}{10.9091}{mi}{103}\RG{86.931}{0.005}{10.9091}{rm}{40}\RG{91.174}{0.005}{10.9091}{mi}{97}\RG{96.945}{0.005}{10.9091}{rm}{41}\RG{104.211}{0.005}{10.9091}{rm}{61}\RG{115.731}{0.005}{10.9091}{rm}{48}\RG{121.186}{0.005}{10.9091}{rm}{46}%
}
\]
这里使用了重根判别：若\allowbreak{}\hskip.1em \RB{100.762}{9.758}{3.208}{%
\RG{0.000}{-0.003}{10.9091}{mi}{70}\RG{8.466}{-0.003}{10.9091}{rm}{40}\RG{12.709}{-0.003}{10.9091}{mi}{120}\RG{18.949}{-0.003}{10.9091}{rm}{41}\RG{26.226}{-0.003}{10.9091}{rm}{61}\RG{37.746}{-0.003}{10.9091}{rm}{40}\RG{41.989}{-0.003}{10.9091}{mi}{120}\RG{50.651}{-0.003}{10.9091}{sy}{0}\RG{61.571}{-0.003}{10.9091}{mi}{97}\RG{67.342}{-0.003}{10.9091}{rm}{41}\RG{71.587}{3.958}{7.9701}{mi}{107}\RG{76.720}{-0.003}{10.9091}{mi}{72}\RG{86.636}{-0.003}{10.9091}{rm}{40}\RG{90.880}{-0.003}{10.9091}{mi}{120}\RG{97.120}{-0.003}{10.9091}{rm}{41}%
}\hskip.1em\allowbreak{} ，其中\allowbreak{}\hskip.1em \RB{25.787}{8.068}{1.801}{%
\RG{0.000}{-0.003}{10.9091}{mi}{107}\RG{8.869}{-0.003}{10.9091}{sy}{21}\RG{20.389}{-0.003}{10.9091}{rm}{50}%
}\hskip.1em\allowbreak{} ，则\allowbreak{}\hskip.1em \RB{83.684}{8.657}{3.208}{%
\RG{0.000}{-0.003}{10.9091}{mi}{70}\RG{8.476}{-0.003}{10.9091}{rm}{40}\RG{12.720}{-0.003}{10.9091}{mi}{97}\RG{18.491}{-0.003}{10.9091}{rm}{41}\RG{25.757}{-0.003}{10.9091}{rm}{61}\RG{37.277}{-0.003}{10.9091}{mi}{70}\RG{45.759}{3.958}{7.9701}{sy}{48}\RG{49.365}{-0.003}{10.9091}{rm}{40}\RG{53.609}{-0.003}{10.9091}{mi}{97}\RG{59.380}{-0.003}{10.9091}{rm}{41}\RG{66.646}{-0.003}{10.9091}{rm}{61}\RG{78.166}{-0.003}{10.9091}{rm}{48}%
}\hskip.1em\allowbreak{} 。

若\allowbreak{}\hskip.1em \RB{41.289}{8.663}{3.202}{%
\RG{0.000}{0.003}{10.9091}{mi}{102}\RG{6.327}{0.003}{10.9091}{rm}{40}\RG{10.571}{0.003}{10.9091}{mi}{97}\RG{16.342}{0.003}{10.9091}{rm}{41}\RG{23.935}{0.003}{10.9091}{rm}{61}\RG{35.771}{0.003}{10.9091}{rm}{48}%
}\hskip.1em\allowbreak{} ，则由第一式得\allowbreak{}\hskip.1em \RB{40.439}{8.663}{3.202}{%
\RG{0.000}{0.003}{10.9091}{mi}{103}\RG{5.476}{0.003}{10.9091}{rm}{40}\RG{9.720}{0.003}{10.9091}{mi}{97}\RG{15.491}{0.003}{10.9091}{rm}{41}\RG{23.084}{0.003}{10.9091}{rm}{61}\RG{34.920}{0.003}{10.9091}{rm}{48}%
}\hskip.1em\allowbreak{} ，于是\allowbreak{}\hskip.1em \RB{25.755}{5.325}{0.617}{%
\RG{0.000}{0.003}{10.9091}{mi}{120}\RG{8.782}{0.003}{10.9091}{sy}{0}\RG{19.822}{0.003}{10.9091}{mi}{97}%
}\hskip.1em\allowbreak{} 同时整除\allowbreak{}\hskip.1em \RB{6.522}{8.194}{2.733}{%
\RG{0.000}{0.003}{10.9091}{mi}{102}%
}\hskip.1em\allowbreak{} 与\allowbreak{}\hskip.1em \RB{5.671}{5.325}{2.733}{%
\RG{0.000}{0.003}{10.9091}{mi}{103}%
}\hskip.1em\allowbreak{} ，与二者互素矛盾。因此

\[
\RB{51.998}{10.104}{3.201}{%
\RG{0.000}{0.004}{10.9091}{mi}{102}\RG{6.316}{0.004}{10.9091}{rm}{40}\RG{10.560}{0.004}{10.9091}{mi}{97}\RG{16.331}{0.004}{10.9091}{rm}{41}\RS{23.607}{0.004}{10.9091}{\ne}\RG{35.128}{0.004}{10.9091}{rm}{48}\RT{40.588}{0.004}{10.9091}{。}%
}
\]
由上述两式，

\[
\RB{245.030}{12.911}{61.470}{%
\RG{173.725}{7.340}{7.9701}{rm}{50}\RG{20.586}{3.963}{7.9701}{rm}{50}\RG{36.360}{3.963}{7.9701}{sy}{48}\RG{54.224}{3.963}{7.9701}{rm}{50}\RG{77.881}{3.963}{7.9701}{sy}{48}\RG{95.745}{3.963}{7.9701}{rm}{50}\RG{151.257}{3.963}{7.9701}{sy}{48}\RG{212.490}{3.963}{7.9701}{rm}{50}\RG{222.810}{3.963}{7.9701}{sy}{48}\RG{240.674}{3.963}{7.9701}{rm}{50}\RG{0.000}{0.002}{10.9091}{mi}{102}\RG{6.327}{0.002}{10.9091}{rm}{40}\RG{10.571}{0.002}{10.9091}{mi}{97}\RG{16.342}{0.002}{10.9091}{rm}{41}\RG{25.429}{8.298}{9.9538}{ex}{0}\RG{30.033}{0.002}{10.9091}{mi}{102}\RG{39.966}{0.002}{10.9091}{rm}{40}\RG{44.210}{0.002}{10.9091}{mi}{97}\RG{49.981}{0.002}{10.9091}{rm}{41}\RG{61.493}{0.002}{10.9091}{rm}{43}\RG{72.402}{0.002}{10.9091}{mi}{103}\RG{81.487}{0.002}{10.9091}{rm}{40}\RG{85.730}{0.002}{10.9091}{mi}{97}\RG{91.501}{0.002}{10.9091}{rm}{41}\RG{100.588}{8.298}{9.9538}{ex}{1}\RG{108.213}{0.002}{10.9091}{rm}{61}\RG{119.733}{8.298}{9.9538}{ex}{0}\RG{124.337}{0.002}{10.9091}{mi}{102}\RG{130.664}{0.002}{10.9091}{rm}{40}\RG{134.908}{0.002}{10.9091}{mi}{97}\RG{140.679}{0.002}{10.9091}{rm}{41}\RG{144.923}{0.002}{10.9091}{mi}{102}\RG{154.863}{0.002}{10.9091}{rm}{40}\RG{159.107}{0.002}{10.9091}{mi}{97}\RG{164.878}{0.002}{10.9091}{rm}{41}\RG{169.122}{8.298}{9.9538}{ex}{1}\RG{180.992}{0.002}{10.9091}{rm}{43}\RG{191.902}{0.002}{10.9091}{mi}{102}\RG{198.229}{0.002}{10.9091}{rm}{40}\RG{202.473}{0.002}{10.9091}{mi}{97}\RG{208.244}{0.002}{10.9091}{rm}{41}\RG{217.334}{0.002}{10.9091}{mi}{103}\RG{226.416}{0.002}{10.9091}{rm}{40}\RG{230.660}{0.002}{10.9091}{mi}{97}\RG{236.430}{0.002}{10.9091}{rm}{41}\RG{139.475}{-16.281}{7.9701}{rm}{50}\RG{149.794}{-16.281}{7.9701}{sy}{48}\RG{167.659}{-16.281}{7.9701}{rm}{50}\RG{206.425}{-16.281}{7.9701}{rm}{50}\RG{216.744}{-16.281}{7.9701}{sy}{48}\RG{234.608}{-16.281}{7.9701}{rm}{50}\RG{108.222}{-20.242}{10.9091}{rm}{61}\RG{119.731}{-20.242}{10.9091}{mi}{103}\RG{125.207}{-20.242}{10.9091}{rm}{40}\RG{129.451}{-20.242}{10.9091}{mi}{97}\RG{135.222}{-20.242}{10.9091}{rm}{41}\RG{144.318}{-20.242}{10.9091}{mi}{103}\RG{153.401}{-20.242}{10.9091}{rm}{40}\RG{157.645}{-20.242}{10.9091}{mi}{97}\RG{163.416}{-20.242}{10.9091}{rm}{41}\RG{174.926}{-20.242}{10.9091}{rm}{43}\RG{185.836}{-20.242}{10.9091}{mi}{102}\RG{192.163}{-20.242}{10.9091}{rm}{40}\RG{196.407}{-20.242}{10.9091}{mi}{97}\RG{202.178}{-20.242}{10.9091}{rm}{41}\RG{211.268}{-20.242}{10.9091}{mi}{103}\RG{220.350}{-20.242}{10.9091}{rm}{40}\RG{224.593}{-20.242}{10.9091}{mi}{97}\RG{230.364}{-20.242}{10.9091}{rm}{41}\RG{144.929}{-36.525}{7.9701}{rm}{50}\RG{182.843}{-36.525}{7.9701}{rm}{50}\RG{197.767}{-36.525}{7.9701}{sy}{48}\RG{215.632}{-36.525}{7.9701}{rm}{50}\RG{108.222}{-40.486}{10.9091}{rm}{61}\RG{119.731}{-32.190}{9.9538}{ex}{0}\RG{124.335}{-40.486}{10.9091}{mi}{102}\RG{130.662}{-40.486}{10.9091}{rm}{40}\RG{134.906}{-40.486}{10.9091}{mi}{97}\RG{140.677}{-40.486}{10.9091}{rm}{41}\RG{152.197}{-40.486}{10.9091}{rm}{43}\RG{163.106}{-40.486}{10.9091}{mi}{103}\RG{168.583}{-40.486}{10.9091}{rm}{40}\RG{172.826}{-40.486}{10.9091}{mi}{97}\RG{178.597}{-40.486}{10.9091}{rm}{41}\RG{187.687}{-32.190}{9.9538}{ex}{1}\RG{192.291}{-40.486}{10.9091}{mi}{103}\RG{201.374}{-40.486}{10.9091}{rm}{40}\RG{205.617}{-40.486}{10.9091}{mi}{97}\RG{211.388}{-40.486}{10.9091}{rm}{41}\RG{108.222}{-60.730}{10.9091}{rm}{61}\RG{119.731}{-60.730}{10.9091}{rm}{48}\RG{125.186}{-60.730}{10.9091}{rm}{46}%
}
\]
由于\allowbreak{}\hskip.1em \RB{40.646}{8.662}{3.203}{%
\RG{0.000}{0.002}{10.9091}{mi}{102}\RG{6.316}{0.002}{10.9091}{rm}{40}\RG{10.560}{0.002}{10.9091}{mi}{97}\RG{16.331}{0.002}{10.9091}{rm}{41}\RS{23.607}{0.002}{10.9091}{\ne}\RG{35.128}{0.002}{10.9091}{rm}{48}%
}\hskip.1em\allowbreak{} ，可得

\[
\RB{93.144}{9.535}{3.202}{%
\RG{6.327}{3.964}{7.9701}{sy}{48}\RG{24.191}{3.964}{7.9701}{rm}{50}\RG{47.847}{3.964}{7.9701}{sy}{48}\RG{65.712}{3.964}{7.9701}{rm}{50}\RG{0.000}{0.003}{10.9091}{mi}{102}\RG{9.933}{0.003}{10.9091}{rm}{40}\RG{14.177}{0.003}{10.9091}{mi}{97}\RG{19.948}{0.003}{10.9091}{rm}{41}\RG{31.459}{0.003}{10.9091}{rm}{43}\RG{42.368}{0.003}{10.9091}{mi}{103}\RG{51.453}{0.003}{10.9091}{rm}{40}\RG{55.697}{0.003}{10.9091}{mi}{97}\RG{61.468}{0.003}{10.9091}{rm}{41}\RG{73.586}{0.003}{10.9091}{rm}{61}\RG{85.095}{0.003}{10.9091}{rm}{48}\RG{90.549}{0.003}{10.9091}{rm}{46}%
}
\]
故\allowbreak{}\hskip.1em \RB{5.933}{5.320}{0.622}{%
\RG{0.000}{-0.002}{10.9091}{mi}{97}%
}\hskip.1em\allowbreak{} 是所述多项式的根。

小结：将重根条件写成函数值与导数值的两条等式，再利用互素性排除公共零点，即可消去非零因子。

\[
\RB{278.451}{8.393}{2.237}{%
\RT{0.000}{0.002}{7.9701}{题}\RT{8.966}{0.002}{7.9701}{目}\RT{17.933}{0.002}{7.9701}{来}\RT{26.899}{0.002}{7.9701}{源}\RT{35.866}{0.002}{7.9701}{：}\RT{44.832}{0.002}{7.9701}{谢}\RT{53.799}{0.002}{7.9701}{启}\RT{62.765}{0.002}{7.9701}{鸿}\RT{71.732}{0.002}{7.9701}{《}\RT{80.698}{0.002}{7.9701}{复}\RT{89.665}{0.002}{7.9701}{旦}\RT{98.631}{0.002}{7.9701}{大}\RT{107.598}{0.002}{7.9701}{学}\RT{116.564}{0.002}{7.9701}{高}\RT{125.531}{0.002}{7.9701}{等}\RT{134.497}{0.002}{7.9701}{代}\RT{143.464}{0.002}{7.9701}{数}\RT{152.430}{0.002}{7.9701}{习}\RT{161.397}{0.002}{7.9701}{题}\RT{170.363}{0.002}{7.9701}{集}\RT{179.329}{0.002}{7.9701}{（}\RT{188.296}{0.002}{7.9701}{2}\RT{192.779}{0.002}{7.9701}{0}\RT{197.262}{0.002}{7.9701}{2}\RT{201.745}{0.002}{7.9701}{3}\RT{209.214}{0.002}{7.9701}{版}\RT{218.180}{0.002}{7.9701}{）}\RT{222.672}{0.002}{7.9701}{》}\RT{227.156}{0.002}{7.9701}{，}\RT{236.122}{0.002}{7.9701}{第}\RT{248.066}{0.002}{7.9701}{6}\RT{252.549}{0.002}{7.9701}{9}\RT{260.018}{0.002}{7.9701}{页}\RT{268.985}{0.002}{7.9701}{。}%
}
\]
% Original page 7.
\section*{2018A14}
设\allowbreak{}\hskip.1em \RB{5.846}{5.327}{2.611}{%
\RG{0.000}{0.005}{10.9091}{mi}{112}%
}\hskip.1em\allowbreak{} 为奇素数。证明：多项式

\[
\RB{223.080}{9.528}{3.208}{%
\RG{74.773}{3.957}{7.9701}{mi}{112}\RG{79.270}{3.957}{7.9701}{sy}{0}\RG{85.212}{3.957}{7.9701}{rm}{50}\RG{142.560}{3.957}{7.9701}{mi}{112}\RG{147.058}{3.957}{7.9701}{sy}{0}\RG{152.999}{3.957}{7.9701}{rm}{51}\RG{0.000}{-0.003}{10.9091}{mi}{102}\RG{6.316}{-0.003}{10.9091}{rm}{40}\RG{10.560}{-0.003}{10.9091}{mi}{120}\RG{16.800}{-0.003}{10.9091}{rm}{41}\RG{24.077}{-0.003}{10.9091}{rm}{61}\RG{35.597}{-0.003}{10.9091}{rm}{40}\RG{39.840}{-0.003}{10.9091}{mi}{112}\RG{47.913}{-0.003}{10.9091}{sy}{0}\RG{58.833}{-0.003}{10.9091}{rm}{49}\RG{64.288}{-0.003}{10.9091}{rm}{41}\RG{68.531}{-0.003}{10.9091}{mi}{120}\RG{92.479}{-0.003}{10.9091}{rm}{43}\RG{103.388}{-0.003}{10.9091}{rm}{40}\RG{107.632}{-0.003}{10.9091}{mi}{112}\RG{115.704}{-0.003}{10.9091}{sy}{0}\RG{126.614}{-0.003}{10.9091}{rm}{50}\RG{132.068}{-0.003}{10.9091}{rm}{41}\RG{136.312}{-0.003}{10.9091}{mi}{120}\RG{160.266}{-0.003}{10.9091}{rm}{43}\RS{171.175}{-0.003}{10.9091}{\cdots}\RG{182.063}{-0.003}{10.9091}{rm}{43}\RG{192.972}{-0.003}{10.9091}{rm}{50}\RG{198.427}{-0.003}{10.9091}{mi}{120}\RG{207.089}{-0.003}{10.9091}{rm}{43}\RG{218.009}{-0.003}{10.9091}{rm}{49}%
}
\]
在有理数域上不可约。

证明：令

\[
\RB{206.736}{9.527}{3.209}{%
\RG{101.969}{3.956}{7.9701}{mi}{112}\RG{106.467}{3.956}{7.9701}{sy}{0}\RG{112.408}{3.956}{7.9701}{rm}{49}\RG{185.808}{3.956}{7.9701}{sy}{48}\RG{0.000}{-0.004}{10.9091}{mi}{104}\RG{6.284}{-0.004}{10.9091}{rm}{40}\RG{10.527}{-0.004}{10.9091}{mi}{120}\RG{16.767}{-0.004}{10.9091}{rm}{41}\RG{24.033}{-0.004}{10.9091}{rm}{61}\RG{35.553}{-0.004}{10.9091}{rm}{49}\RG{43.429}{-0.004}{10.9091}{rm}{43}\RG{54.350}{-0.004}{10.9091}{mi}{120}\RG{63.011}{-0.004}{10.9091}{rm}{43}\RS{73.921}{-0.004}{10.9091}{\cdots}\RG{84.808}{-0.004}{10.9091}{rm}{43}\RG{95.728}{-0.004}{10.9091}{mi}{120}\RG{117.251}{-0.004}{10.9091}{rm}{44}\RG{143.913}{-0.004}{10.9091}{mi}{102}\RG{150.240}{-0.004}{10.9091}{rm}{40}\RG{154.484}{-0.004}{10.9091}{mi}{120}\RG{160.724}{-0.004}{10.9091}{rm}{41}\RG{168.001}{-0.004}{10.9091}{rm}{61}\RG{179.521}{-0.004}{10.9091}{mi}{104}\RG{189.414}{-0.004}{10.9091}{rm}{40}\RG{193.657}{-0.004}{10.9091}{mi}{120}\RG{199.897}{-0.004}{10.9091}{rm}{41}\RG{204.141}{-0.004}{10.9091}{rm}{46}%
}
\]
由二项式定理，

\[
\RB{191.107}{17.899}{13.855}{%
\RG{134.223}{14.031}{7.9701}{mi}{112}\RG{89.065}{11.351}{7.9701}{mi}{112}\RG{55.546}{7.390}{10.9091}{rm}{40}\RG{59.790}{7.390}{10.9091}{mi}{120}\RG{68.452}{7.390}{10.9091}{rm}{43}\RG{79.361}{7.390}{10.9091}{rm}{49}\RG{84.816}{7.390}{10.9091}{rm}{41}\RG{96.485}{7.390}{10.9091}{sy}{0}\RG{107.394}{7.390}{10.9091}{rm}{49}\RG{154.195}{7.390}{10.9091}{mi}{112}\RG{174.116}{3.965}{7.9701}{mi}{106}\RG{177.728}{3.965}{7.9701}{sy}{0}\RG{183.669}{3.965}{7.9701}{rm}{49}\RG{0.000}{0.004}{10.9091}{mi}{104}\RG{6.284}{0.004}{10.9091}{rm}{40}\RG{10.527}{0.004}{10.9091}{mi}{120}\RG{19.189}{0.004}{10.9091}{rm}{43}\RG{30.098}{0.004}{10.9091}{rm}{49}\RG{35.553}{0.004}{10.9091}{rm}{41}\RG{42.829}{0.004}{10.9091}{rm}{61}\RG{117.077}{0.004}{10.9091}{rm}{61}\RG{128.595}{10.368}{10.9092}{ex}{88}\RG{146.166}{15.431}{8.7012}{ex}{32}\RG{159.846}{15.431}{8.7012}{ex}{33}\RG{167.876}{0.004}{10.9091}{mi}{120}\RG{188.513}{0.004}{10.9091}{rm}{46}\RG{81.079}{-7.480}{10.9091}{mi}{120}\RG{154.703}{-7.480}{10.9091}{mi}{106}\RG{129.523}{-11.797}{7.9701}{mi}{106}\RG{133.135}{-11.797}{7.9701}{rm}{61}\RG{139.076}{-11.797}{7.9701}{rm}{49}\RR{55.546}{2.732}{57.304}{0.436}%
}
\]
对两边求导，得

\[
\RB{187.787}{17.890}{13.863}{%
\RG{100.807}{14.022}{7.9701}{mi}{112}\RG{150.876}{7.382}{10.9091}{mi}{112}\RG{170.796}{3.956}{7.9701}{mi}{106}\RG{174.408}{3.956}{7.9701}{sy}{0}\RG{180.350}{3.956}{7.9701}{rm}{50}\RG{0.000}{-0.004}{10.9091}{mi}{70}\RG{8.476}{-0.004}{10.9091}{rm}{40}\RG{12.720}{-0.004}{10.9091}{mi}{120}\RG{18.960}{-0.004}{10.9091}{rm}{41}\RG{26.226}{-0.004}{10.9091}{rm}{58}\RG{29.258}{-0.004}{10.9091}{rm}{61}\RG{40.779}{-0.004}{10.9091}{mi}{102}\RG{47.106}{-0.004}{10.9091}{rm}{40}\RG{51.350}{-0.004}{10.9091}{mi}{120}\RG{60.011}{-0.004}{10.9091}{rm}{43}\RG{70.932}{-0.004}{10.9091}{rm}{49}\RG{76.386}{-0.004}{10.9091}{rm}{41}\RG{83.652}{-0.004}{10.9091}{rm}{61}\RG{95.179}{10.360}{10.9092}{ex}{88}\RG{110.932}{-0.004}{10.9091}{rm}{40}\RG{115.176}{-0.004}{10.9091}{mi}{106}\RG{122.234}{-0.004}{10.9091}{sy}{0}\RG{133.143}{-0.004}{10.9091}{rm}{49}\RG{138.598}{-0.004}{10.9091}{rm}{41}\RG{142.841}{15.423}{8.7012}{ex}{32}\RG{156.526}{15.423}{8.7012}{ex}{33}\RG{164.555}{-0.004}{10.9091}{mi}{120}\RG{185.192}{-0.004}{10.9091}{rm}{46}\RG{151.383}{-7.489}{10.9091}{mi}{106}\RG{96.107}{-11.805}{7.9701}{mi}{106}\RG{99.719}{-11.805}{7.9701}{rm}{61}\RG{105.660}{-11.805}{7.9701}{rm}{50}%
}
\]
其首项系数为\allowbreak{}\hskip.1em \RB{23.846}{7.764}{2.617}{%
\RG{0.000}{-0.001}{10.9091}{mi}{112}\RG{7.964}{-0.001}{10.9091}{sy}{0}\RG{18.775}{-0.001}{10.9091}{rm}{49}%
}\hskip.1em\allowbreak{} ，不被\allowbreak{}\hskip.1em \RB{5.846}{5.321}{2.617}{%
\RG{0.000}{-0.001}{10.9091}{mi}{112}%
}\hskip.1em\allowbreak{} 整除。对\allowbreak{}\hskip.1em \RB{63.031}{7.764}{2.737}{%
\RG{0.000}{-0.001}{10.9091}{rm}{50}\RG{8.476}{-0.001}{10.9091}{sy}{20}\RG{19.997}{-0.001}{10.9091}{mi}{106}\RG{27.666}{-0.001}{10.9091}{sy}{20}\RG{39.186}{-0.001}{10.9091}{mi}{112}\RG{47.160}{-0.001}{10.9091}{sy}{0}\RG{57.961}{-0.001}{10.9091}{rm}{49}%
}\hskip.1em\allowbreak{} ，\allowbreak{}\hskip.1em \RB{12.940}{9.194}{5.824}{%
\RG{0.000}{8.295}{9.9538}{ex}{0}\RG{4.604}{4.843}{7.9701}{mi}{112}\RG{5.047}{-3.766}{7.9701}{mi}{106}\RG{9.102}{8.295}{9.9538}{ex}{1}%
}\hskip.1em\allowbreak{} 的分子含因子\allowbreak{}\hskip.1em \RB{5.846}{5.321}{2.617}{%
\RG{0.000}{-0.001}{10.9091}{mi}{112}%
}\hskip.1em\allowbreak{} ，而分母\allowbreak{}\hskip.1em \RB{7.220}{8.310}{2.737}{%
\RG{0.000}{-0.001}{10.9091}{mi}{106}\RG{4.625}{-0.001}{10.9091}{rm}{33}%
}\hskip.1em\allowbreak{} 不含因子\allowbreak{}\hskip.1em \RB{5.846}{5.321}{2.617}{%
\RG{0.000}{-0.001}{10.9091}{mi}{112}%
}\hskip.1em\allowbreak{} ，

\noindent \RB{267.705}{27.740}{10.824}{%
\RT{0.000}{17.640}{10.9091}{故}\RG{251.431}{7.382}{10.9091}{mi}{112}\RG{196.742}{-0.004}{10.9091}{mi}{112}\RG{205.414}{-0.004}{10.9091}{sy}{106}\RG{211.480}{-0.004}{10.9091}{rm}{40}\RG{215.724}{-0.004}{10.9091}{mi}{106}\RG{222.782}{-0.004}{10.9091}{sy}{0}\RG{233.702}{-0.004}{10.9091}{rm}{49}\RG{239.157}{-0.004}{10.9091}{rm}{41}\RG{243.400}{15.423}{8.7012}{ex}{32}\RG{257.081}{15.423}{8.7012}{ex}{33}\RG{265.110}{-0.004}{10.9091}{rm}{46}\RG{251.938}{-7.489}{10.9091}{mi}{106}%
}

常数项为

\[
\RB{116.911}{16.276}{10.819}{%
\RG{44.996}{7.387}{10.9091}{mi}{112}\RG{74.421}{7.387}{10.9091}{mi}{112}\RG{80.061}{7.387}{10.9091}{rm}{40}\RG{84.304}{7.387}{10.9091}{mi}{112}\RG{92.388}{7.387}{10.9091}{sy}{0}\RG{103.297}{7.387}{10.9091}{rm}{49}\RG{108.752}{7.387}{10.9091}{rm}{41}\RG{0.000}{0.001}{10.9091}{mi}{70}\RG{8.476}{0.001}{10.9091}{rm}{40}\RG{12.720}{0.001}{10.9091}{rm}{48}\RG{18.175}{0.001}{10.9091}{rm}{41}\RG{25.440}{0.001}{10.9091}{rm}{61}\RG{36.960}{15.428}{8.7012}{ex}{32}\RG{50.647}{15.428}{8.7012}{ex}{33}\RG{61.698}{0.001}{10.9091}{rm}{61}\RG{114.196}{0.001}{10.9091}{rm}{44}\RG{45.094}{-7.483}{10.9091}{rm}{50}\RG{90.983}{-7.483}{10.9091}{rm}{50}\RR{74.421}{2.729}{38.579}{0.436}%
}
\]
由\allowbreak{}\hskip.1em \RB{5.846}{5.321}{2.617}{%
\RG{0.000}{-0.001}{10.9091}{mi}{112}%
}\hskip.1em\allowbreak{} 为奇素数，\allowbreak{}\hskip.1em \RB{37.439}{8.681}{3.228}{%
\RG{0.000}{-0.001}{10.9091}{mi}{112}\RG{9.120}{-0.001}{10.9091}{sy}{106}\RG{15.622}{-0.001}{10.9091}{mi}{70}\RG{24.098}{-0.001}{10.9091}{rm}{40}\RG{28.342}{-0.001}{10.9091}{rm}{48}\RG{33.797}{-0.001}{10.9091}{rm}{41}%
}\hskip.1em\allowbreak{} 而\allowbreak{}\hskip.1em \RB{43.486}{9.531}{3.228}{%
\RG{0.000}{-0.001}{10.9091}{mi}{112}\RG{5.651}{3.960}{7.9701}{rm}{50}\RS{13.967}{-0.001}{10.9091}{\nmid}\RG{21.669}{-0.001}{10.9091}{mi}{70}\RG{30.145}{-0.001}{10.9091}{rm}{40}\RG{34.389}{-0.001}{10.9091}{rm}{48}\RG{39.844}{-0.001}{10.9091}{rm}{41}%
}\hskip.1em\allowbreak{} 。因此\allowbreak{}\hskip.1em \RB{8.669}{7.917}{0.501}{%
\RG{0.000}{-0.001}{10.9091}{mi}{70}%
}\hskip.1em\allowbreak{} 满足素数\allowbreak{}\hskip.1em \RB{5.846}{5.321}{2.617}{%
\RG{0.000}{-0.001}{10.9091}{mi}{112}%
}\hskip.1em\allowbreak{} 对应的Eisenstein判别条件，故在\allowbreak{}\hskip.1em \RB{7.166}{8.190}{2.617}{%
\RG{0.000}{-0.001}{10.9091}{bb}{81}%
}\hskip.1em\allowbreak{}上不可约。

为说明判别依据，注意\allowbreak{}\hskip.1em \RB{85.341}{8.656}{3.209}{%
\RG{0.000}{-0.004}{10.9091}{mi}{70}\RG{8.476}{-0.004}{10.9091}{rm}{40}\RG{12.720}{-0.004}{10.9091}{sy}{0}\RG{21.207}{-0.004}{10.9091}{rm}{49}\RG{26.662}{-0.004}{10.9091}{rm}{41}\RG{33.928}{-0.004}{10.9091}{rm}{61}\RG{45.448}{-0.004}{10.9091}{mi}{102}\RG{51.775}{-0.004}{10.9091}{rm}{40}\RG{56.019}{-0.004}{10.9091}{rm}{48}\RG{61.473}{-0.004}{10.9091}{rm}{41}\RG{68.750}{-0.004}{10.9091}{rm}{61}\RG{80.270}{-0.004}{10.9091}{rm}{49}%
}\hskip.1em\allowbreak{} ，故\allowbreak{}\hskip.1em \RB{8.669}{7.914}{0.504}{%
\RG{0.000}{-0.004}{10.9091}{mi}{70}%
}\hskip.1em\allowbreak{} 是本原整系数多项式。若它在\allowbreak{}\hskip.1em \RB{7.166}{8.187}{2.620}{%
\RG{0.000}{-0.004}{10.9091}{bb}{81}%
}\hskip.1em\allowbreak{} 上可约，由Gauss引理可写为

\[
\RB{233.479}{9.291}{3.229}{%
\RG{0.000}{-0.002}{10.9091}{mi}{70}\RG{11.498}{-0.002}{10.9091}{rm}{61}\RG{23.018}{-0.002}{10.9091}{mi}{71}\RG{31.593}{-0.002}{10.9091}{mi}{72}\RG{41.509}{-0.002}{10.9091}{rm}{44}\RG{68.182}{-0.002}{10.9091}{mi}{71}\RG{76.757}{-0.002}{10.9091}{rm}{44}\RG{81.601}{-0.002}{10.9091}{mi}{72}\RG{94.550}{-0.002}{10.9091}{sy}{50}\RG{104.859}{-0.002}{10.9091}{bb}{90}\RG{112.136}{-0.002}{10.9091}{rm}{91}\RG{115.234}{-0.002}{10.9091}{mi}{120}\RG{121.474}{-0.002}{10.9091}{rm}{93}\RG{124.507}{-0.002}{10.9091}{rm}{44}\RT{151.181}{-0.002}{10.9091}{d}\RT{157.247}{-0.002}{10.9091}{e}\RT{162.090}{-0.002}{10.9091}{g}\RG{169.363}{-0.002}{10.9091}{mi}{71}\RG{177.938}{-0.002}{10.9091}{rm}{44}\RT{182.789}{-0.002}{10.9091}{d}\RT{188.855}{-0.002}{10.9091}{e}\RT{193.699}{-0.002}{10.9091}{g}\RG{200.972}{-0.002}{10.9091}{mi}{72}\RG{213.910}{-0.002}{10.9091}{sy}{21}\RG{225.430}{-0.002}{10.9091}{rm}{49}\RG{230.885}{-0.002}{10.9091}{rm}{46}%
}
\]
Gauss引理在此可简证如下：两个本原整系数多项式的乘积仍本原，否则模某个整除全部乘积系数的素数约化，将得到两个非零多项式之积为零的矛盾；清除有理分解的分母，再提取各因子的系数公因子，即得到上述整数分解：

将\allowbreak{}\hskip.1em \RB{41.704}{8.193}{0.738}{%
\RG{0.000}{0.002}{10.9091}{mi}{70}\RG{11.498}{0.002}{10.9091}{rm}{61}\RG{23.018}{0.002}{10.9091}{mi}{71}\RG{31.593}{0.002}{10.9091}{mi}{72}%
}\hskip.1em\allowbreak{} 模\allowbreak{}\hskip.1em \RB{5.846}{5.324}{2.614}{%
\RG{0.000}{0.002}{10.9091}{mi}{112}%
}\hskip.1em\allowbreak{} 约化，有

\[
\RB{85.234}{9.534}{2.613}{%
\RG{67.358}{3.963}{7.9701}{mi}{112}\RG{71.855}{3.963}{7.9701}{sy}{0}\RG{77.797}{3.963}{7.9701}{rm}{50}\RG{0.000}{0.003}{10.9091}{mi}{71}\RG{10.392}{0.003}{10.9091}{mi}{72}\RG{23.330}{0.003}{10.9091}{rm}{61}\RG{34.857}{0.003}{10.9091}{mi}{112}\RG{42.930}{0.003}{10.9091}{sy}{0}\RG{53.839}{0.003}{10.9091}{rm}{49}\RG{61.116}{0.003}{10.9091}{mi}{120}\RG{82.640}{0.003}{10.9091}{rm}{46}\RR{0.000}{9.221}{8.574}{0.436}\RR{10.392}{8.981}{9.917}{0.436}\RR{34.857}{8.795}{24.442}{0.436}%
}
\]
\noindent \RB{465.396}{26.358}{3.228}{%
\RT{0.000}{16.258}{10.9091}{由}\RT{11.051}{16.258}{10.9091}{于}\RG{25.806}{16.258}{10.9091}{mi}{112}\RT{35.297}{16.258}{10.9091}{不}\RT{46.348}{16.258}{10.9091}{整}\RT{57.410}{16.258}{10.9091}{除}\RG{72.161}{16.258}{10.9091}{mi}{70}\RT{84.477}{16.258}{10.9091}{的}\RT{95.528}{16.258}{10.9091}{首}\RT{106.590}{16.258}{10.9091}{项}\RT{117.641}{16.258}{10.9091}{系}\RT{128.703}{16.258}{10.9091}{数}\RT{139.612}{16.258}{10.9091}{，}\RT{150.663}{16.258}{10.9091}{两}\RT{161.725}{16.258}{10.9091}{因}\RT{172.787}{16.258}{10.9091}{子}\RT{183.838}{16.258}{10.9091}{约}\RT{194.900}{16.258}{10.9091}{化}\RT{205.951}{16.258}{10.9091}{后}\RT{217.012}{16.258}{10.9091}{次}\RT{228.074}{16.258}{10.9091}{数}\RT{239.125}{16.258}{10.9091}{不}\RT{250.187}{16.258}{10.9091}{变}\RT{261.097}{16.258}{10.9091}{；}\RT{272.148}{16.258}{10.9091}{而}\RT{283.209}{16.258}{10.9091}{单}\RT{294.271}{16.258}{10.9091}{项}\RT{305.322}{16.258}{10.9091}{式}\RT{316.384}{16.258}{10.9091}{的}\RT{327.435}{16.258}{10.9091}{因}\RT{338.497}{16.258}{10.9091}{子}\RT{349.559}{16.258}{10.9091}{仍}\RT{360.610}{16.258}{10.9091}{是}\RT{371.672}{16.258}{10.9091}{单}\RT{382.723}{16.258}{10.9091}{项}\RT{393.785}{16.258}{10.9091}{式}\RT{404.694}{16.258}{10.9091}{，}\RT{415.756}{16.258}{10.9091}{所}\RT{426.807}{16.258}{10.9091}{以}\RG{441.562}{16.258}{10.9091}{mi}{71}\RT{453.977}{16.258}{10.9091}{与}\RG{243.085}{3.959}{7.9701}{rm}{50}\RG{0.000}{-0.001}{10.9091}{mi}{72}\RT{13.548}{-0.001}{10.9091}{均}\RT{24.457}{-0.001}{10.9091}{被}\RG{38.998}{-0.001}{10.9091}{mi}{120}\RT{48.870}{-0.001}{10.9091}{整}\RT{59.780}{-0.001}{10.9091}{除}\RT{70.689}{-0.001}{10.9091}{。}\RT{81.598}{-0.001}{10.9091}{于}\RT{92.507}{-0.001}{10.9091}{是}\RG{107.043}{-0.001}{10.9091}{mi}{112}\RG{115.716}{-0.001}{10.9091}{sy}{106}\RG{121.781}{-0.001}{10.9091}{mi}{71}\RG{130.356}{-0.001}{10.9091}{rm}{40}\RG{134.599}{-0.001}{10.9091}{rm}{48}\RG{140.054}{-0.001}{10.9091}{rm}{41}\RT{147.936}{-0.001}{10.9091}{且}\RG{162.477}{-0.001}{10.9091}{mi}{112}\RG{171.150}{-0.001}{10.9091}{sy}{106}\RG{177.216}{-0.001}{10.9091}{mi}{72}\RG{187.132}{-0.001}{10.9091}{rm}{40}\RG{191.376}{-0.001}{10.9091}{rm}{48}\RG{196.830}{-0.001}{10.9091}{rm}{41}\RT{201.081}{-0.001}{10.9091}{，}\RT{211.990}{-0.001}{10.9091}{从}\RT{222.899}{-0.001}{10.9091}{而}\RG{237.434}{-0.001}{10.9091}{mi}{112}\RG{250.959}{-0.001}{10.9091}{sy}{106}\RG{257.014}{-0.001}{10.9091}{mi}{70}\RG{265.490}{-0.001}{10.9091}{rm}{40}\RG{269.734}{-0.001}{10.9091}{rm}{48}\RG{275.188}{-0.001}{10.9091}{rm}{41}\RT{279.440}{-0.001}{10.9091}{，}\RT{290.349}{-0.001}{10.9091}{矛}\RT{301.259}{-0.001}{10.9091}{盾}\RT{312.168}{-0.001}{10.9091}{。}\RT{323.077}{-0.001}{10.9091}{这}\RT{333.986}{-0.001}{10.9091}{就}\RT{344.895}{-0.001}{10.9091}{给}\RT{355.805}{-0.001}{10.9091}{出}\RT{366.714}{-0.001}{10.9091}{了}\RT{377.623}{-0.001}{10.9091}{本}\RT{388.532}{-0.001}{10.9091}{题}\RT{399.441}{-0.001}{10.9091}{所}\RT{410.350}{-0.001}{10.9091}{用}\RT{421.260}{-0.001}{10.9091}{判}\RT{432.169}{-0.001}{10.9091}{别}\RT{443.078}{-0.001}{10.9091}{法}\RT{453.987}{-0.001}{10.9091}{的}\RR{441.562}{25.476}{8.575}{0.436}\RR{0.000}{8.977}{9.916}{0.436}%
}

证明：

最后，代换\allowbreak{}\hskip.1em \RB{47.366}{7.762}{1.409}{%
\RG{0.000}{-0.004}{10.9091}{mi}{120}\RG{9.262}{-0.004}{10.9091}{sy}{55}\RG{22.953}{-0.004}{10.9091}{mi}{120}\RG{31.495}{-0.004}{10.9091}{rm}{43}\RG{42.295}{-0.004}{10.9091}{rm}{49}%
}\hskip.1em\allowbreak{} 是\allowbreak{}\hskip.1em \RB{18.904}{8.678}{3.231}{%
\RG{0.000}{-0.004}{10.9091}{bb}{81}\RG{7.276}{-0.004}{10.9091}{rm}{91}\RG{10.375}{-0.004}{10.9091}{mi}{120}\RG{16.615}{-0.004}{10.9091}{rm}{93}%
}\hskip.1em\allowbreak{} 的自同构，其逆为\allowbreak{}\hskip.1em \RB{47.355}{7.762}{0.624}{%
\RG{0.000}{-0.004}{10.9091}{mi}{120}\RG{9.262}{-0.004}{10.9091}{sy}{55}\RG{22.953}{-0.004}{10.9091}{mi}{120}\RG{31.495}{-0.004}{10.9091}{sy}{0}\RG{42.284}{-0.004}{10.9091}{rm}{49}%
}\hskip.1em\allowbreak{} ，故保持可约性。因此\allowbreak{}\hskip.1em \RB{6.522}{8.187}{2.740}{%
\RG{0.000}{-0.004}{10.9091}{mi}{102}%
}\hskip.1em\allowbreak{} 在\allowbreak{}\hskip.1em \RB{7.166}{8.187}{2.620}{%
\RG{0.000}{-0.004}{10.9091}{bb}{81}%
}\hskip.1em\allowbreak{} 上不可约。

小结：把几何级数多项式的导数平移到\allowbreak{}\hskip.1em \RB{24.642}{7.770}{1.401}{%
\RG{0.000}{0.004}{10.9091}{mi}{120}\RG{8.662}{0.004}{10.9091}{rm}{43}\RG{19.571}{0.004}{10.9091}{rm}{49}%
}\hskip.1em\allowbreak{} ，即可得到满足Eisenstein判别法的二项式系数。

\[
\RB{278.451}{8.393}{2.237}{%
\RT{0.000}{0.002}{7.9701}{题}\RT{8.966}{0.002}{7.9701}{目}\RT{17.933}{0.002}{7.9701}{来}\RT{26.899}{0.002}{7.9701}{源}\RT{35.866}{0.002}{7.9701}{：}\RT{44.832}{0.002}{7.9701}{谢}\RT{53.799}{0.002}{7.9701}{启}\RT{62.765}{0.002}{7.9701}{鸿}\RT{71.732}{0.002}{7.9701}{《}\RT{80.698}{0.002}{7.9701}{复}\RT{89.665}{0.002}{7.9701}{旦}\RT{98.631}{0.002}{7.9701}{大}\RT{107.598}{0.002}{7.9701}{学}\RT{116.564}{0.002}{7.9701}{高}\RT{125.531}{0.002}{7.9701}{等}\RT{134.497}{0.002}{7.9701}{代}\RT{143.464}{0.002}{7.9701}{数}\RT{152.430}{0.002}{7.9701}{习}\RT{161.397}{0.002}{7.9701}{题}\RT{170.363}{0.002}{7.9701}{集}\RT{179.329}{0.002}{7.9701}{（}\RT{188.296}{0.002}{7.9701}{2}\RT{192.779}{0.002}{7.9701}{0}\RT{197.262}{0.002}{7.9701}{2}\RT{201.745}{0.002}{7.9701}{3}\RT{209.214}{0.002}{7.9701}{版}\RT{218.180}{0.002}{7.9701}{）}\RT{222.672}{0.002}{7.9701}{》}\RT{227.156}{0.002}{7.9701}{，}\RT{236.122}{0.002}{7.9701}{第}\RT{248.066}{0.002}{7.9701}{6}\RT{252.549}{0.002}{7.9701}{9}\RT{260.018}{0.002}{7.9701}{页}\RT{268.985}{0.002}{7.9701}{。}%
}
\]

\end{document}
